\documentclass{llncs}
\usepackage{xspace}
\usepackage{epsfig}
\usepackage{amsmath, amssymb}
\usepackage{algorithmicx,algorithm}
\usepackage[noend]{algpseudocode}
\usepackage{pifont}
\newcommand\comment[1]{}
\usepackage{listings}
\usepackage{tkz-graph}
\usepackage{graphics}
%\usepackage[small,it]{caption}
\usetikzlibrary{arrows,shapes}
\lstset{language=C}


% Macros
\newcommand\limply{\rightarrow}
\newcommand\type[1]{{\tt #1}}
\newcommand\nin{\noindent}
\newcommand\true{\mbox{{\sc true}}\xspace}
\newcommand\false{\mbox{{\sc false}}\xspace}
% \newcommand\comment[1]{}
% \newcommand\note[2]{\begin{itemize} \item[$\Diamond$ \emph{Side-Note $\rangle$}] {\bf #1: }#2 $\Diamond$\end{itemize}}
\newcommand\brackets[1]{\vspace{0.5 cm}\begin{center}#1\end{center}\vspace{0.3 cm}}
\newcommand\alg[1]{{\sc #1}}
\newcommand\margparams[1]{\marginpar{\fbox{#1}}}
\newcommand\textfirst[1]{\emph{#1}\index{#1}}
%\newcommand\first[2]{\margin{#2}\label{page:#1}}
% we need a separate macro for each chapter because we keep the name of the figure equal to the name of the .fig file, and we cannot use ':' in the file name
\newcommand\graph[2]{\begin{figure} \begin{center}\scalebox{0.7}{\input{#1.pstex_t}}\caption{\label{fig-eq:#1}#2}\end{center}\end{figure}}
\newcommand\graphsa[2]{\begin{figure} \begin{center}\scalebox{0.7}{\input{#1.pstex_t}}\caption{\label{fig-sa:#1}#2}\end{center}\end{figure}}
\newcommand\term[1]{\emph{#1}\index{#1}}
\newcommand\function[1]{{\sc #1}}
\newcommand\tool[1]{{\sc #1}}
\newcommand\SEM[1]{\mathopen{{[}\kern-.25ex{[}}\;#1\;\mathclose{]}\kern-.25ex{]}}
\newcommand\rulename[1]{{\sc #1}\index{inference rules!#1}\xspace}
\newcommand\infrule[3]{\displaystyle{\frac{#1}{#2}}\ \ (\mbox{\rulename{#3}})}
%\newcommand\corn{\ \lower 5pt\hbox{\hskip 4pt\vrule width 3pt height 1pt depth
%  1pt\hbox{\vrule width 2pt height 4.5pt depth 1pt}}}        % end of def

\newcommand\Corn{\hfill $\diamond$}

\newcommand\corn{\hfill\hbox{\hskip 4pt
                \vrule width 3pt height 1pt depth 1pt
                \hbox{\vrule width 2pt height 3.5pt depth 1pt}}}


% eq specific macros



\newcommand\E{\mbox{\tiny E}} % for superscript E
\newcommand\eq{\varphi^{\E}} % Equality formula (after reduction)
\newcommand\peq{{\psi^{\E}}}  % Equality formula (after reduction)
\newcommand\equf{\varphi^{\mbox{\tiny UF}}}% Equality formula with Uninterpreted Functions (before reduction)
\newcommand\eqfc{consistent^{\E}} % Functional consistency constraints
\newcommand\eqflat{flat^{\E}} % same as \equf, when ufs are replaced with variables
\newcommand\eqF{\varphi^{\E}_\star} % same as \equf, when ufs are replaced withthe \F variables
\newcommand\bool{{\mathcal B}} % Propositional formula
\newcommand\enc{\bool_{enc}} % Boolean encoding
\newcommand\transitivity{\bool_{trans}} % Transitivity constraints due to Boolean encoding
\newcommand\transsparse{\bool_{sparse}} % Transitivity constraints due to the sparse method
\newcommand\transopt{\bool_{rtc}} % Transitivity constraints due to the reduced transitivity constraints method
\newcommand\alphaopt{\alpha_{rtc}} % assignment due to the reduced transitivity constraints method
\newcommand\alphasparse{\alpha_{sparse}} % assignment due to the sparse method
\newcommand\eqgraph{G^{\E}}
\newcommand\eqp{=^*} % Equality path
\newcommand\neqp{\neq^*} % disequality path
\newcommand\eqsp{=^*_s}
\newcommand\neqsp{\neq^*_s}
\newcommand\size[1]{\vert #1\vert}
\newcommand\Ene{E_{\neq}} % set of \neq edges
\newcommand\Eeq{E_=} %  set of = edges
\newcommand\Geq{\eqgraph_=} %   dashed sub-graph (representing \Eeq)
\newcommand\Gne{\eqgraph_{\neq}} % solid sub-graph (representing \Ene)
\newcommand\Gin{G_{\emph{inc}}}
\newcommand\cha{\emph{char}}
\newcommand\F{F^\star}
\newcommand\G{G^\star}
% from cav01
\newtheorem{Yrule}{Rule}
\newtheorem{YTrule}{(Suggested) Rule}
\newcommand{\simple}[1]{{\it simp}({#1})}
% uf specific macros
\newcommand\Fcase[1]{$\left( \mbox{\begin{tabular}{lll}#1 \end{tabular}} \right)$}
\newcommand\Fcaseb[1]{$\left( \mbox{\begin{tabular}{lllll}#1 \end{tabular}} \right)$}

%\renewcommand\vec[1]{\overline{#1}}
\renewcommand\term[1]{term(#1)}
\newcommand\is{\mathbb{L_{UP}}} % inference system with UP
\newcommand\lmt{\mathbb{L_{RE}}} % inference system with MT
\newcommand\codesize{\footnotesize}
\newcommand\recur{{\sc (proc-p-eq)}\xspace}
\newcommand\recurs{{\sc (proc-p-eq$_s$)}\xspace}
\newcommand\recurn{proc-p-eq}
\newcommand\mtermg{{\sc (m-term-pe)}\xspace}
\newcommand\mtermd{{\sc (m-term)}\xspace}
\newcommand\mtermdp{{\sc (m-term$^+$)}\xspace}
\newcommand\termd{{\sc (term)}\xspace}
\newcommand\callc{{\sc call\mbox{-}contains}\xspace}
\newcommand\mapf{map_{\cal F}}
\newcommand\mapg{map_g}
\newcommand\mapt{map_t}
\newcommand\mapm{map_m}
\newcommand\mapl[1]{$\langle #1 \rangle \in map$} % inside the first alg.
\newcommand\map[1]{$\langle #1 \rangle \in \mapf$}
\newcommand\upbody[1]{{#1}^{UF}}
\newcommand\upbodiesand[1]{$\upbody{#1}$ and $\upbody{#1'}$}
\newcommand\upproc{\mbox{UF}}
\newcommand\upprocl[1]{\textbf{\upproc(#1)}}
\newcommand\upprocsand[1]{$\upproc(#1)$ and $\upproc(#1')$}
\newcommand\upprocs[1]{$\upproc(#1)$, $\upproc(#1')$}
\newcommand\ufpdef[1]{PV_f[#1]}
\newcommand\ufpdefp[1]{PV_{f'}[#1]}

% for the eval expression tree example:
\newcommand\isleaf{\mbox{isleaf}}
\newcommand\isplus{\mbox{isplus}}
\newcommand\ismult{\mbox{ismult}}
\newcommand\mkplus{\mbox{mkplus}}
\newcommand\mkmult{\mbox{mkmult}}
\newcommand\leftson{\mbox{leftson}}
\newcommand\rightson{\mbox{rightson}}
\newcommand\doplus{\mbox{doplus}}
\newcommand\domult{\mbox{domult}}

\newcommand\from{\mathit{from}}
\newcommand\RT{RT\xspace}
\newcommand\RV{RV\xspace}
\newcommand\Tree{Stack-level tree}
\newcommand\tree{stack-level tree}
\newcommand\trees{stack-level trees}



\newcommand\cb{check-block\xspace}
\newcommand\cbr{check-block$^r$\xspace}
\newcommand\cbs{check-blocks\xspace}
%\newcommand\uf{uninterpreted function\xspace}
\newcommand\ufs{uninterpreted functions\xspace}
\newcommand\UFs{Uninterpreted Functions\xspace}
\newcommand\UF{Uninterpreted Function\xspace}
\newcommand\rsp{related subprogram}
\newcommand\rsps{related subprograms}
\newcommand\lrf{loop recursive function}
\newcommand\lrfs{loop recursive functions}
\newcommand\ufstr{UF-struct}
\newcommand\ufstrs{UF-structs}
\newcommand\ufarr{UF-array}
\newcommand\ufarrs{UF-arrays}
\newcommand\ub{unbounded structure}
\newcommand\ubs{unbounded structures}
\newcommand\ua{uninterpreted array}
\newcommand\uas{uninterpreted arrays}

\newcommand\difflabel{{``Failed''}\xspace}

\newcommand\equivscclabel{{``Conditional-Equivalent''}\xspace}

\newcommand\call{\mbox{call }}
\newcommand\body{\mbox{ body}}
\newcommand\nextup{\alg{next-unmarked-pair()}\xspace}
\newcommand\Checkr{Check$^r$}
\newcommand\Long[1]{} % redefine for long version
\newcommand\pe{p\mbox{-}equiv\xspace}
\newcommand\requiv{call\mbox{-}equiv}
\newcommand\rimplic{\xrightarrow{reach}}


\newcommand\mtbody[1]{{#1}^{UF}}
\newcommand\mtproc{\mbox{UF}}
\newcommand\mtprocl[1]{{\sc uf}\_{#1}}
\newcommand\mtufprocl[1]{{\sc peuf}\_{#1}}
\newcommand\pneuf[1]{{\sc uf}(#1, $\vec{in}$)}
\newcommand\pneufp[1]{{\sc uf'}(#1, $\vec{in'}$)}
\newcommand\peuf[1]{{\sc peuf}(#1, $\vec{in}$)}
\newcommand\peufp[1]{{\sc peuf'}(#1, $\vec{in'}$)}
\newcommand\uf[1]{UF_{#1}}
\newcommand\ufp[1]{\ufpdef{\uf{#1}}}
\newcommand\ufpp[1]{\ufpdefp{\uf{#1}}}

\newcommand\pneufname{{\sc uf}\xspace}
\newcommand\pneufpname{{\sc uf'}\xspace}
\newcommand\peufname{{\sc peuf}\xspace}
\newcommand\peufpname{{\sc peuf'}\xspace}

\newcommand\emulatesame{\xspace{\sc ($^=$)}\xspace}
\newcommand\emulatedifferent{\xspace{\sc ($^{\neq}$)}\xspace}

\newcommand\mtname{reach\mbox{-}equivalent\xspace}
\newcommand\mt{m\mbox{-}term\xspace}
\newcommand\comtag[1]{$#1$, $#1'$}
\newcommand\comandtag[1]{$#1$ and $#1'$}
\newcommand\upcomtag[1]{\upproc($#1$), \upproc($#1'$)}
\newcommand\mscccovered{$covered$\xspace}
\newcommand\ProveMT{\alg{ProveMT}\xspace}
\newcommand\CheckCallEquiv{\alg{CallEquiv}\xspace}
\newcommand\abort{{\bf abort}}
\newcommand\mtlabel{$m\_term$\xspace}
\newcommand\equivlabel{$partially\_equiv$\xspace}
\newcommand\coverlabel{$mscc\_covered$\xspace}
\newcommand\pair[1]{\langle#1\rangle}
\newcommand\pairtag[1]{$\langle$\comtag{#1}$\rangle$}
\newcommand\pairbodies[1]{$\pair{\upbody{#1}$, $\upbody{#1'}}$}
\newcommand\enforcement{{\sc (Enforce)}\xspace}
\newcommand\enforcementnew{{\sc (Enforce-new)}\xspace}
\newcommand\attachmentProveMT{Fig.~\ref{alg:attachment}}

\newcommand\requivpair[1]{\requiv{(#1, #1'})}
\newcommand\mtpair[1]{\mt{(#1, #1'})}
\newcommand\calls[2]{#1\xrightarrow{calls}#2}
\newcommand\pairwtag[1]{\pair{#1,#1'}}
\newcommand\callspair[2]{\calls{\pairwtag{#1}}{\pairwtag{#2}}}

\newcommand\full[1]{#1}
\newcommand\notfull[1]{}
\newcommand\cg{\sc{cg}}
\newcommand\student{ }
\addtolength{\intextsep}{-0.5 cm}
\addtolength{\abovecaptionskip}{-0.3 cm}
\begin{document}
\title{Proving Mutual Termination of Programs \full{(full version)}}


\author{
Dima Elenbogen$^1$ \quad
Shmuel Katz$^1$\quad
Ofer Strichman$^2$
}
\institute{
$^1$ CS, Technion, Haifa, Israel. \email{\{katz,edima\}@cs.technion.ac.il}  \\
$^2$ Information Systems Engineering, IE, Technion, Haifa, Israel\ \email{ofers@ie.technion.ac.il}\\
}

\newcommand\bpar[1]{{\bf #1} }

\maketitle

\begin{abstract}
Two programs are said to be \emph{mutually terminating} if they terminate on
exactly the same inputs. We suggest a proof rule that uses a mapping between the functions of the two programs for proving mutual
termination of functions \comtag{f}. The rule's premise requires proving that given the same
arbitrary input $\vec{in}$, $f(\vec{in})$ and $f'(\vec{in})$ call mapped functions with the same
arguments. A variant of this rule with a weaker premise allows to prove
termination of one of the programs if the other is known to terminate for all
inputs. We present an algorithm for decomposing the verification problem of
whole programs to that of proving mutual termination of individual functions,
based on our suggested rules.
\end{abstract}

\section{Introduction} \label{sec:intro}

Whereas termination of a single program has been widely studied~(e.g.,
\cite{Fl67,CPR05,BMS05,CPR11}) for several decades by now, with the focus
being, especially in the last few years, on automating such proofs,
little attention has been paid to the related problem of proving that
two similar programs (e.g., two consecutive versions of the same
program) terminate on exactly the same inputs.  Ideally one should
focus on the former problem, but this is not always possible either
because the automatic techniques are inherently incomplete, or because
\emph{by design} the program does not terminate on all inputs. In such
cases there is value in solving the latter problem, because
developers may wish to know that none of their changes affect the
termination behavior of their program. Moreover, the problem and
solution thereof can be defined in the granularity of functions rather
than whole programs; in this case the developer may benefit even more from
a detailed list of pairs of functions that terminate on exactly the
same set of inputs. Those pairs that are not on the list can help
detecting termination errors.

Our focus is on successive, closely related versions of a program because it
both reflects a realistic problem of developers, and offers opportunities for
decomposition and abstraction that are not possible with the single-program
termination problem.  This problem, which was initially proposed in~\cite{GS08}
and coined \emph{mutual termination}, can easily be proven undecidable as can
be seen via a simple reduction from the halting problem. We argue, however,
that in many cases it is easier to solve automatically, because unlike termination proofs for a single program, it does not rely on proving that
the sequence of states in the programs' computations can be mapped into a well-founded
set. Rather it can be proven by showing that the loops and recursive functions
have the same set of function calls given the same inputs, which is relatively
easier to prove automatically. In Sec.~\ref{sec:rule}, for example, we show how to prove mutual termination of two versions of the famous Collatz's $3x + 1$ problem~\cite{G81}; whereas proving termination of this program is open for many decades, proving mutual termination with respect to another version is simple.

Our suggested method for decomposing the proof is most valuable when the two
input programs $P$ and $P'$ are relatively similar in structure. In fact, its
complexity is dominated by the \emph{difference} between the programs, rather
than by their absolute size. It begins by heuristically building a (possibly
partial) map between the functions of $P$ and $P'$. It then progresses
bottom-up on the two call graphs, and each time proves the mutual termination
of a pair of functions in the map, while abstracting their callees. The
generated verification conditions are in the form of assertions about `flat' programs (i.e., without
loops and recursive calls), which are proportional in size to the two compared
functions. It then discharges these verification conditions with a bounded
model-checker (CBMC~\cite{CK03} in our case). Each such program has the same
structure: it calls the two compared functions sequentially with the same
nondeterministic input, records all subsequent function calls and their
arguments, and asserts in the end that they have an equivalent set of function
calls. According to our proof rule, the validity of this assertion is
sufficient for establishing their mutual termination.


The algorithm is rather involved because it has to deal with cases in which the
call graphs of $P$ and $P'$ are not isomorphic (this leads to unmapped
functions), with mutually recursive functions, and with cases in which the
proof of mutual termination for the callees has failed. It also improves completeness by utilizing extra
knowledge that we may give to it on the \emph{partial equivalence} of the
callees, where two functions are said to be partially equivalent if given the
same inputs they terminate with the same outputs, or at least one of them does
not terminate.
%If we know that two mapped callees are partially equivalent, we
%abstract them with the same uninterpreted function,
%which increases our chance
%to prove mutual termination.
Partial equivalence was studied in~\cite{GS08,GS09} and is implemented in \tool{rvt}~\cite{GS09} and Microsoft's
\tool{SymDiff}~\cite{KLR10}. We also implemented our algorithm in \tool{rvt}, which enables us to gain this information in a preprocessing step.


To summarize our contributions in this paper, we present a) a proof rule for inferring mutual termination of recursive (and mutually-recursive) functions at the leaves of their respective call graphs,
%, which requires a weaker premise than the one in~\cite{GS10}
 b) an extension of the first rule that applies also to internal nodes in the call graphs, and c) a proof rule for inferring \emph{termination} (not mutual termination) in case the other function is known to be terminating. More importantly, we show how these rules can be applied to whole programs via a bottom-up decomposition algorithm, and report on a prototype implementation of this algorithm -- the first to deal with the mutual termination problem.

%In the next section we give a formal definition of our problem and a description of the preprocessing that we %apply to programs in order to be able to use our proof rules and decomposition algorithm.
%Section~\ref{sec:rule} describe a proof rule for proving mutual termination of functions in mutual recursion. Section~\ref{sec:decomposition} suggests a method for applying the rule for whole programs, based on a bottom-up traversal of the two call graphs. In Sec.~\ref{sec:term} we consider a related problem: \emph{assuming} that $P$ terminates, prove that $P'$ terminates. Experiments and conclusions are summarized in Sec.~\ref{sec:experiments}.

\section{Preliminaries} \label{sec:base}
%From here on we will study the problem of mutual termination of
%functions. Programs can be seen as a special case, namely the 'main' function.
Our goal is to prove mutual termination of pairs of functions in programs that are assumed to be deterministic (i.e.,
single threaded and no internal nondeterminism). Formally:

\begin{definition}[Mutual termination of functions]
\label{def:mt_prog} Two functions $f$ and $f'$  are \emph{mutually
terminating} if and only if they terminate upon exactly the same inputs.
\end{definition}
By \emph{input} we mean both the function parameters and the
global data it accesses, including the heap. 
%
Denote by $\mt(f,f')$ the fact that $f$ and $f'$ mutually terminate.
%
\subsubsection{Preprocessing and mapping}
As a preprocessing step, all loops are extracted to external recursive
functions, as shown in~\cite{G08}. After this step nontermination can only
arise from recursion. In addition, all global variables that are read by a
function are added to the end of its formal parameter list, and the
calling sites are changed accordingly. This is not essential for the
proof, but simplifies the presentation. It should be noted that this step in
itself is impossible in general programs that access the heap, because it is
undecidable whether there exists an input to a function that causes the
function to read a particular variable. Our only way out of this problem is to
point out that it is easy to overapproximate this information (in the worst
case just take the whole list of global variables) and to state that, based on
our experience with a multitude of real programs, it is rather easy to compute
this information precisely or slightly overapproximate it with static analysis
techniques such as alias analysis. Indeed, the same exact problem exists in
\tool{rvt} and \tool{SymDiff} for the case of partial equivalence, and there,
as in our case, overapproximation can only hinder completeness, not soundness.
In general we will not elaborate on issues arising from aliasing because these
are not unique to mutual termination, and are dealt with in~\cite{GS09,KLR10}.

As a second step, a map $\mapf$ between the functions of $P$ and $P'$ is
derived. For functions $f \in P$ and $f' \in P'$ it is possible that
$\pair{f,f'} \in \mapf$ only if $f$ and $f'$ have the same prototype, i.e., the
same list of formal input parameter types. We emphasize that the output of the two
functions need not be compatible (e.g., $f$ can update more global variables
than $f'$).

\subsubsection{Definitions and notations}

\begin{itemize}
\item \emph{Function isolation}. With each function $g$, we associate an
    uninterpreted function $\uf{g}$ such that $g$ and $\uf{g}$ have the
    same prototype and return type\footnote{This definition generalized naturally to cases in which $g$ has multiple outputs owing to global data and arguments passed by reference.}. The following definition will be used for specifying which functions are associated with the \emph{same} uninterpreted function:    
    
\begin{definition}[Partial equivalence of functions]
\label{def:pe_prog} Two functions $f$ and $f'$  are \emph{partially
equivalent} if any two terminating executions of $f$ and $f'$ starting from the
same inputs, return the same value.
\end{definition}
Denote by $\pe(f,f')$ the fact that $f$ and $f'$ are partially equivalent. We enforce that
\begin{equation}\label{eq:enforcement}
\uf{g} = \uf{g'} \Rightarrow \left(\pair{g, g'} \in\mapf \wedge \pe(g, g')\right)  \mbox{(enforce)}
\end{equation}
i.e., we associate $g$ and $g'$ with the same
uninterpreted function only if $\pair{g,g'} \in \mapf$, and $g,g'$ were proven to be partially equivalent. The list of pairs of functions that are proven to be partially equivalent is assumed to be an input to the mutual termination algorithm.
%
We now define:
\begin{equation}\label{eq:isolation}
\mtbody{f} \doteq f[g(expr_{in}) \leftarrow \uf{g}(expr_{in}) \mid g \mbox{ is called in $f$}]\;,
\end{equation}
%
where $expr_{in}$ is the expression(s) denoting actual parameter(s) with which $g$ is called.
$ \mtbody{f}$ is called the \emph{isolated} version of $f$. By construction
it has no loops or function calls, except for calls to uninterpreted functions.

The definition of $\mtbody{f}$ requires all function calls to be replaced
with uninterpreted functions. A useful relaxation of this requirement,
which we will later use, is that it can inline non-recursive functions.
Clearly the result is still nonrecursive. Therefore, we still refer to this
as an isolated version of $f$.


\item \emph{Call equivalence}.
%
%The following definition is necessary for the proof rule:
%
%
%\begin{definition}[Call-containment]
%\label{def:rimplic_func} The calls of a function $f$ are said to
%\emph{contain} those of $f'$ if the existence of a call to $g'(\vec{in})$
%in $f'$ implies that there is a call to $g(\vec{in})$ in every possible
%computation of $f$, given that \pairtag{g} $\in\mapf$, and that $f,f'$ were
%called with equal inputs.
%\end{definition}
%

%*** new

Let $calls(f(\vec{in}))$, where $\vec{in}$ is a vector of actual values, denote the set of function call instructions (i.e., a function name and the actual parameter values) invoked in the body of $f()$ during the execution of $f(\vec{in})$. Note that $calls(f(\vec{in}))$ does not include calls from descendant functions and hence also not from recursive calls.

We can now define:
\begin{definition}[Call-equivalence of functions]\label{def:re_func}
$f$ and $f'$ are \emph{call-equivalent} if and only if
\[ \forall \text{\pairtag{g}} \in \mapf, in_f, in_g.\ g'(in_g) \in calls(f'(in_f))
\Leftrightarrow g(in_g) \in calls (f(in_f))\;. \]
\end{definition}
%
%Denote by $\callc(f,f')$ that the calls of $f$ contain those of $f'$.
%
%\begin{definition}[Call-equivalence of functions]
%\label{def:re_func} Two functions \comandtag{f} are said to be
%\emph{call-equivalent} if $\callc(f,f') \land \callc(f',f)$.
%%, when called with the same input arguments, for every call that $f$ makes
%%to a function $g$, there is a call by $f'$ to $g'$ (and vice-versa) for
%%every possible computation such that \pairtag{g} $\in\mapf$, and \comtag{g}
%%are called with the same arguments.
%\end{definition}
Denote by $\requiv(f,f')$ the fact that $f$ and $f'$ are call-equivalent.
%
Note that it is decidable whether $f^{UF}$ and $f'^{UF}$ are call-equivalent, because
these are functions without loops and function calls, as explained above.

\end{itemize}

\section{A proof rule}
\label{sec:rule}
Given a call graph of a general program, a corresponding DAG may be built by collapsing each maximal strongly connected
component (MSCC) into a single node. Nodes that are not part of any cycle in
the call graph (corresponding to non-recursive functions) are called
\emph{trivial} MSCCs in the DAG. Other MSCCs correspond to either simple or
mutually recursive function(s).

Consider, then, two nontrivial MSCCs \comtag{m} in $P$ and $P'$, respectively,
that are \emph{leaves} in the MSCC DAGs.
%(the rule in~\cite{GS08} only refers to this case). !!! add note to related work
Also  assume, for now, that $\mapf$ provides a bijective mapping between the functions in \comandtag{m}. Denote by
\[ \mapf(m) = \{\pair{f,f'}  \mid \pair{f,f'} \in \mapf, f \in m\}\;.\]
%
Our goal is to prove mutual termination of each of the pairs in $\mapf(m)$. The
following proof rule gives us a way to do it by proving
call-equivalence of each of these pairs:
\begin{center}
\fbox{
\begin{minipage}{\linewidth}
\begin{equation} \label{eq:new}
\frac
{\forall \pair{f, f'} \in\mapf(m).\ \requiv(\upbody{f},\upbody{f'})}
{\forall \pair{f, f'} \in\mapf(m).\ \mt(f, f')}
\mbox{\quad \mtermd}
\end{equation}
\end{minipage}
}
\end{center}
%
Later on in this section we generalize this rule so it also applies to non-leaf MSCCs,
and outline the soundness proof.

%
Readers familiar with~\cite{GS08} may wonder about the relation of this rule to the one suggested there. The answer is that (\ref{eq:new}) is stronger (i.e., it has a weaker premise) because it does not require the functions to be partially equivalent.

Note that the abstraction of calls with uninterpreted functions is the source of incompleteness. $\requiv(\upbody{f}, \upbody{f'})$ may fail, but the counterexample may rely on values returned by an uninterpreted function that are different than what the corresponding concrete function would have returned if called with the same parameters. Furthermore, it is possible that the concrete function and its counterpart on the other side do not terminate, but their abstractions terminate and are followed by different function calls on the two sides, which leads to call equivalence not being true.



\subsubsection{Checking the premise}
We check the premise of~(\ref{eq:new}) by building a loop- and recursion-free program for each pair of functions that we want to prove call equivalent. Here we describe the construction informally, and only for the case of simple recursion at the leaf functions. We will consider the general case in a more formal way in Sec.~\ref{sec:decomposition}.

Let $f,f'$ be simple recursive functions that only call themselves. We associate a set of call instructions with each called function (this set represents $calls(f(\vec{in}))$. For example, in $f$ only $f$ itself is called, and hence we maintain a set of call instructions to $f$. We then build a program with the following structure: {\tt main} assigns equal nondeterministic values to the inputs of $f$ and $f'$. It then calls an implementation of $\mtbody{f}$ and $\mtbody{f'}$, and finally asserts that the sets of call instructions are equal. The example below (hopefully) clarifies this construction.
% The %implementation of $\mtbody(f)$ is simply $f$ itself where the recursive call is replaced with a call to a new function that a) models an uninterpreted function and b) collect call








\begin{figure}[h!]
\begin{tabular}{p{6 cm}p{6 cm}} \hline
\begin{algorithmic}
\Function {$f$}{{\bf int} $a$}
\State {\bf int} $even := 0, ret := 0$;
\If {$a > 1$}
\If {$\lnot(a\ \%\ 2)$} \Comment{even}
    \State $a := a/2$;
    \State $even = 1$;
    \Else \ $a := 3a + 1$;
\EndIf
\State $ret := even + f(a)$;
\EndIf
\Return $ret$;
\EndFunction
\end{algorithmic}
&
\begin{algorithmic}
\Function{$f'$}{{\bf int} a'}
\State {\bf int} $t', odd' := 0, ret' := 0$;
\If {$a' \leq 1$} \Return $ret'$; \EndIf
\State 	$t' := a' / 2$;
\If {$a' \% 2$}   \Comment{odd}
\State		$a' := 6 t' + 4$;
\State		$odd' := 1$;
\Else \ $a' := t'$;	
\EndIf
\State $ret' := odd' + f'(a')$;
\State \Return $ret'$;
\EndFunction
\end{algorithmic} \\[-6pt] \hline
\end{tabular}
\caption{Two variations on the Collatz (``$3x + 1$") function that are mutually terminating. $f$ ($f'$) returns the total number of times the function was called with an even (odd) number. Note than when $a'$ is odd, $a' / 2 = (a'-1)/2$, and hence $6(a'/2) + 4 = 3a' + 1$.}\label{fig:collatz}
\end{figure}



\begin{example} \label{ex:collatz} Consider the two variants of the Collatz (``$3x + 1$") program~\cite{G81} in Fig.~\ref{fig:collatz}\footnote{In the pseudocode we use the convention by which $\%$ is the modulo operator. \Long{, `:=' is an assignment and `=' is equality.}}, which return different values (see explanation in the caption of the figure).
%\footnote{the code uses C-like notation. Hence `$>>$' is `shift-right' and `$\&$' is `bit-wise and'},
The Collatz program is a famous open problem in termination: no one knows whether
it terminates for all (unbounded) integers. On
the other hand proving mutual termination of the two variants given here is
easy. The comparison is not fair, however, because our decision procedure assumes finite types: we target C programs.
But as we show in \notfull{the full version of this article~\cite{EKS12}}\full{Appendix~\ref{sec:infinite}}, it is solvable even when the input parameter is an unbounded integer, using a decision procedure for linear arithmetic.


\begin{figure}[h!]
\begin{tabular}{p{6 cm}p{6 cm}} \hline
\begin{algorithmic}
\Function {$f^{UF}$}{{\bf int} $a$}
\State {\bf int} $even := 0, ret := 0$;
\If {$a > 1$}
\If {$\lnot(a\ \%\ 2)$} \Comment{even}
    \State $a := a/2$;
    \State $even = 1$;
    \Else \ $a := 3a + 1$;
\EndIf
\State $ret := even$ + \alg{UF}($f,a$);
\EndIf
\Return $ret$;
\EndFunction
\end{algorithmic}
&
\begin{algorithmic}
\Function {$f'^{UF}$}{{\bf int} $a'$}
\State {\bf int} $t', odd' := 0, ret' := 0$;
\If {$a' \leq 1$} \Return $ret'$; \EndIf
\State 	$t' := a' / 2$;
\If {$a' \% 2$}   \Comment{odd}
\State		$a' := 6 t' + 4$;
\State		$odd' := 1$;
\Else \ $a' := t'$;	
\EndIf
\State $ret' := odd'$ + \alg{UF}($f',a'$);
\State \Return $ret'$;
\EndFunction
\end{algorithmic}
\\
%
\multicolumn{2}{c}{\hspace{0.3 cm}
\begin{minipage}{11.6 cm}
\begin{algorithmic}  % implementation of UFf and UFf'
\Function{uf}{function index $g$, input parameters $\vec{in}$}
  \If {$\vec{in} \in params[g]$} \Return the output of the earlier call \pneuf{$g$};
  \EndIf
  \State $params[g] := params[g]\; \bigcup\; \vec{in}$;
  \State \Return a nondeterministic value; % nondet();
  \EndFunction\newline
\end{algorithmic}
\end{minipage}
} \\
\multicolumn{2}{c}{\hspace{0.3 cm}
\begin{minipage}{11.6 cm}
\begin{algorithmic}
\Function{main}{}
\State $\vec{in} = nondet()$; $f^{UF}(\vec{in})$; $f'^{UF}(\vec{in})$;
\State assert($params[f] = params[f']$); \Comment{checks call equivalence}
\EndFunction
\end{algorithmic}
\end{minipage}
}
\\ \mbox{} \\ \hline
\end{tabular}
%
\caption{The flat program that we generate and then verify its assertion, given the two functions of Fig.~\ref{fig:collatz} .} \label{fig:collatz-flat}
\end{figure}



% Note than when $a'$ is odd, $a' >> 1 = (a'-1)/2$, and hence $6 \cdot (a'>>1) + 4 = 3a' + 1$.
%Note that we made the two functions return \emph{different} values, just in order to emphasize that mutual termination can be proved for functions that are not partially equivalent.
The definitions of $f^{UF}$, $f'^{UF}$ appear at the top part of Fig.~\ref{fig:collatz-flat}. The middle part of the same figure shows an implementation {\tt UF} of the uninterpreted functions. It receives a function index (abusing notation for simplicity, we assume here that a function name represents also a unique index) and the actual parameters. Note that it records the set of call instructions in the array {\tt params}.






Note that in this case $f,f'$ are not partially equivalent, and therefore according to $(\ref{eq:enforcement})$ we replace the recursive calls with different uninterpreted functions. Indeed, we call \alg{uf} above with two different function indices ($f$ and $f'$), which means that on equal inputs they do not necessarily return the same nondeterministic value.
We defer the presentation of the case in which the functions are known to be partially equivalent to Sec.~\ref{sec:decomposition}.
%
\qed
\end{example}

What if the mapping is not bijective, or if some of the pairs cannot be proven to be mutually terminating? It is not hard to see that it is sufficient to prove mutual termination of pairs of functions that together intersect all cycles in \comtag{m}, whereas the other functions are inlined. The same observation was made with regard to proving partial equivalence in~\cite{GS11}.
This observation can be used to improve completeness: even when there is no bijective mapping or when it is impossible to prove mutual termination for all pairs in \comtag{m}, it is still sometimes possible to prove it for some of the pairs. The algorithm that we describe in Sec.~\ref{sec:decomposition} uses this observation.

\comment{Moreover, \mtermd only refers to leaf MSCCs. The algorithm by which whole programs are checked for mutual termination has never been given explicitly in previous publications nor was it implemented. Questions like how to deal with non-leaf MSCCs and how to progress
once the descendants are not equivalent were left open, and are newly treated here (in Sec.~\ref{sec:decomposition} we will give an
explicit decomposition algorithm, which also uses the observation above).
}

\Long{We continue in the next section by generalizing \mtermd to the case in which functions are called outside the MSCCs.}

\subsubsection{Generalization} We now generalize \mtermd to the case that \comtag{m} are not leaf MSCCs. This means that there is a set of functions $C(m)$ outside of $m$ that are called by functions in $m$. $C(m')$ is defined similarly with respect to $m'$. The
premise now requires that these functions are mutually-terminating:

\begin{center}
\fbox{
\begin{minipage}{\linewidth}
\begin{equation} \label{eq:new1}
\frac
{\begin{array}{l}
(\forall \pair{g, g'} \in \mapf.\ (g \in C(m) \land g' \in C(m')) \rightarrow \mt(g,g')) \land \\
(\forall \pair{f, f'} \in\mapf(m).\ \requiv(\upbody{f},\upbody{f'}))
\end{array}
}
{\forall \pair{f, f'} \in\mapf(m).\ \mt(f, f')}
\mbox{\ \mtermdp} \;.
\end{equation}
\end{minipage}
}
\end{center}
%
Recall that (\ref{eq:isolation}) prescribes that calls to functions in $C(m)$ and $C(m')$ are replaced with uninterpreted functions in $\upbody{f}$, $\upbody{f'}$.

A full soundness proof of the generalized rule appears in Appendix~\ref{sec:proof1}, whereas
here we only sketch its steps. The proof begins by showing that the premise implies $\forall \pair{f, f'} \in\mapf(m).\ \requiv(f,f')$.
%the validity of:
%
%$ \big(\forall \pair{f, f'} \in\mapf(m).\ \requiv(\upbody{f},\upbody{f'})\big)
%\rightarrow \forall \pair{f, f'} \in\mapf(m).\ \requiv(f,f') \;. $
\Long{
\[ \begin{array}{l} \big(\forall \pair{f, f'} \in\mapf(m).\ \requiv(\upbody{f},\upbody{f'})\big)
\rightarrow \\
\forall \pair{f, f'} \in\mapf(m).\ \requiv(f,f')\;. \end{array} \]}
%
Now, falsely assume that there is a pair \pairtag{f}$\in \mapf(m)$ that is not mutually
terminating whereas the premise holds.
For some value $in$, suppose that it is $f(in)$ that terminates, while $f'(in)$
does not. The infinite call stack of $f'(in)$ must contain a call, say from
$h'(in_1)$ to $g'(in_2)$, whereas $h(in_1)$ does not call $g(in_2)$ in the call
stack of $f(in)$ (assuming $\{$\pairtag{g}, \pairtag{h}$\} \subseteq \mapf$).
This contradicts our premise that \pairtag{h} are call-equivalent. The argument
is a little more involved when there are multiple calls to the same function, and when there are calls to functions in $C(m), C(m')$,
but we leave such subtleties to Appendix~\ref{sec:proof1}.

\section{A decomposition algorithm}
\label{sec:decomposition}

In this section we present an algorithm for proving mutual termination of full programs.
As mentioned in Sec.~\ref{sec:rule}, the call graph of a program can be viewed
as a DAG where the nodes correspond to MSCCs. After building a mapping between the MSCCs of the two call graphs, the algorithm traverses the DAG
bottom-up. For each mapped pair of MSCCs $m,m'$, it attempts to prove the mutual termination of their mapped functions, based on \mtermdp.

The algorithm is inspired by a similar algorithm for verification of \emph{partial
equivalence}, which is described in a technical report~\cite{GS11}
\Long{\footnote{\emph{Note to reviewers}: this report is also currently under review
for a journal~\cite{GS10}}}. The algorithm here is more involved, however,
because it handles differently cases in which the checked functions are also
partially equivalent (recall that this information, i.e., which functions are
known to be partially equivalent, is part of the input to the algorithm).
Furthermore, the algorithm in~\cite{GS11} is described with a non-deterministic
step, and here we suggest a method for determinizing it.

The preprocessing and mapping is as in Sec.~\ref{sec:base}. Hence the program is loop-free, globals accessed by a function are sent instead as additional inputs, and there is a (possibly partial) mapping $\mapf$ between the functions of $P$ and $P'$.

%Note that wrong mapping does not affect soundness: it is used for generating the verification conditions. Hence, wrong mapping can only cause a proof failure.

%For a simpler description, the algorithm assumes that each function receives only one parameter. But this is not a limitation: we can consider a function that receives several parameters as if it gets a tuple, i.e., one complex parameter. Thus, the aforementioned expectation of the same input parameters amounts to an expectation of equivalent tuple structures among the members of each function pair.

\subsection{The algorithm}
\Long{
Table~\ref{tab:markings} contains the labels we use in the decomposition
algorithm (Alg.~\ref{alg:ProveMT}).

\begin{table}[here]
\begin{tabular}{||l|l|p{6.2cm}||} \hline \label{table:markings}
\textbf{Predicate label}	& \textbf{What labeled}  & \textbf{Meaning} \\ \hline
\mtlabel      & \pairtag{f}$ \in\mapf$ & The mapped functions \comandtag{f} are mutually terminating \\ \hline
\equivlabel   & \pairtag{f}$ \in\mapf$ & The mapped functions \comandtag{f} are partially equivalent \\ \hline
\mscccovered  & \pairtag{m}$ \in\mapm$ & The MSCC pair has been processed. \\ \hline
\end{tabular}
\caption{Predicate labels used in the decomposition algorithm.}
\label{tab:markings}
\end{table}
}
The input to Alg.~\ref{alg:ProveMT} is \comtag{P}, a (possibly partial) mapping $\mapf$ between their functions, and (implicitly) those paired functions that are known to be partially equivalent. Its output is a set of function pairs that are marked as \mtlabel, indicating it succeeded to prove their mutual termination based on \mtermdp. We now describe the three functions used by this algorithm.


\subsubsection{\ProveMT.} \label{sec:ProveMT}
This entry function traverses the call
graphs of $P,P'$ bottom-up, each time focusing on a pair of MSCCs. In
line~\ref{step:inline} it inlines all nonrecursive functions that are not
mapped. In line $\ref{step:rename}$ it uses renaming to resolve possible name
collisions between the globals of the two input programs. The next line builds
the MSCC DAGs $MD$ and $MD'$ from the call graphs, as explained in
Sec.\ref{sec:rule}. Line~\ref{step:generate_map} attempts to build a bijective
mapping $\mapm$ between the nodes of $MD$ and $MD'$ which is consistent with
$map_f$. Namely, if $\pair{m,m'} \in \mapm$, $f$ is a function in $m$, and
$\pair{f,f'} \in \mapf$, then $f'$ is a function in $m'$ (and vice-versa). If
such a mapping is impossible, the algorithm aborts. In practice one may run the
algorithm bottom-up until reaching nonmapped MSCCs, but we omit this option
here for brevity.

The bottom-up traversal starts in line~\ref{step:while}. Initially all MSCCs
are unmarked. The algorithm searches for a next unmarked pair \pairtag{m} of
MSCCs such that all its children pairs are marked. If \comtag{m} are trivial
(see Sec.~\ref{sec:rule} for a definition), then line~\ref{step:Check} simply
checks the call-equivalence of the function pair \pairtag{f} that constitutes
\pairtag{m}, and marks them accordingly in line~\ref{step:markTrivial}. 
Note that even if the descendants of \comtag{m} are mutually-terminating, \comtag{m}
are not necessarily so, because they may call their descendants with different
parameters. Also note that if this check fails, we continue to check their ancestors, (in contrast to the case of non-trivial MSCCs listed next), because even if \pairtag{f} are not
mutually terminating for every input, their callers may still be (they can be mutually
terminating in the context of their callers). We can check this by inlining
them, which is only possible because they are not recursive.



\noindent
\begin{algorithm}
\begin{minipage}{\linewidth}
\begin{algorithmic}[1]
\Function{ProveMT}{Programs \comtag{P}, map between functions $\mapf$}
\State \label{step:inline} Inline non-recursive non-mapped functions;
\State \label{step:rename} Solve name collisions in global identifiers of $P,P'$ by
renaming.
\State \label{step:generate} Generate MSCC DAGs \comtag{MD} from the
call graphs of $P,P'$;

\State If possible, generate a bijective map
\label{step:generate_map} $\mapm$ between the nodes of $MD$ \newline \mbox{~~~~~~} and $MD'$ that is consistent with $\mapf$;
\Long{\footnote{It is desirable but not necessary to add pairs of trivial nodes to $\mapm$.}}
Otherwise \abort.

\While {$\exists \pair{m,m'} \in \mapm$
          not marked \mscccovered but its children are,} \label{step:while}
   \State \label{step:choose} Choose such a pair $\pair{m,m'} \in \mapm$ and mark it \mscccovered
   \If {$m,m'$ are trivial} \label{step:m1}
     \State Let $f,f'$ be the functions in $m,m'$, respectively;
     \If {\CheckCallEquiv(\alg{isolate}$(f,f',\emptyset)$)} \label{step:Check}
     mark $f,f'$ as \mtlabel; \label{step:markTrivial}
     \EndIf
   \Else
    \label{step:select} Select non-deterministically
         $S \subseteq \{\pair{f,f'}\mid \pair{f,f'}\in\mapf(m)\}$ \newline
          \mbox{~~~~~~~~~~~~~~~~} that intersect every cycle in \comandtag{m};
      \For {each $\pair{f,f'} \in S$}\label{step:forall}
        \If  {$\lnot$\CheckCallEquiv(\alg{isolate}$(f,f',S)$)} \label{step:abort}{mark ancestors as \mscccovered;}
        \EndIf
      \EndFor
      \For {each $\pair{f,f'} \in S$} \label{step:forall2}
         mark $f,f'$ as \mtlabel; \label{step:markNonTrivial}
      \EndFor
   \EndIf
\EndWhile
\EndFunction%\newline

\vspace{0.25 cm}


\Function{isolate}{functions \comtag{f}, function pairs $S$} \Comment{Builds
$\mtbody{f}, \mtbody{f'}$}
  \For {each $\{\pair{g,g'}\in \mapf \mid g,g'$ are reachable from $f,f'$\}}\label{step:forcheckCallequiv}
    \If {$\pair{g,g'} \in S$ or $\pair{g,g'}$ is marked \mtlabel} \label{step:substitute}
      \State \label{step:subst_uf}Replace calls to $g(\vec{in})$, $g'(\vec{in'})$ with calls to \pneuf{g}, \pneufp{g'}, resp.;
    \Else { inline $g,g'$ in their callers; } \label{step:isinline}
    \EndIf
  \EndFor
  \State \Return $\pair{f,f'}$;
\EndFunction
\vspace{0.25 cm}

\Function {CallEquiv}{A pair of isolated functions $\pair{\mtbody{f}, \mtbody{f'}}$}
\label{step:checkCallEquiv}
  \State \label{step:gen_main} Let $\delta$ denote the program:
  \[
  \begin{array}{l}
 \rhd \mbox{ here add the definitions of {\sc uf()} and {\sc uf'()} (see Fig.~\ref{fig:ufs1}).}\\
 \vec{in} := nondet(); \mtbody{f}(\vec{in}); \mtbody{f'}(\vec{in}); \\
% \mbox{assert}(params(${\sc UF}$) \subseteq params(${\sc UF'}$));\\
 \forall \pair{g,g'} \mbox{ called in } f,f'.\ \mbox{assert}(params[g] \subseteq params[g']);\\
 \end{array}
 \]
  \State \Return CBMC($\delta$); \label{step:cbmc}
\EndFunction

\end{algorithmic}
\end{minipage}
\caption{Pseudo-code for a bottom-up decomposition algorithm for proving that pairs of functions mutually terminate.}
\label{alg:ProveMT}
\end{algorithm}


Next, consider the case that the selected \comtag{m} in line~\ref{step:choose}
are not trivial. In line~\ref{step:select} the algorithm chooses
non-deterministically a subset $S$ of pairs from $\mapf(m)$ that intersects all
the cycles in \comandtag{m}\Long{\footnote{In graph-theoretic terms, the functions in
$S$ constitute a \emph{feedback vertex set} of both $m$ and $m'$\Long{~\cite{E79}}.}}.
This guarantees that we can always inline the functions in \comtag{m}
that are not in $S$. Determinization of this step will be considered in
subsection~\ref{sec:determinization}. If \CheckCallEquiv returns \true for all
the function pairs in $S$, then all those pairs are labeled as \mtlabel in
line~\ref{step:forall2}. Otherwise it abandons the attempt to prove their ancestors in line~\ref{step:abort}: it cannot prove
that mapped functions in $\pair{m,m'}$ are mutually
terminating, nor can it inline these functions in their callers, so we cannot check all its ancestors.

Regardless of whether \pairtag{m} are trivial, they get marked as \coverlabel
in line~\ref{step:choose}, and the loop in \ProveMT  continues to another
pair.




\comment{
\begin{algorithm}[t]
\begin{minipage}{\linewidth}
\begin{algorithmic}[1]

\Function{isolate}{functions \comtag{f}, function pairs $S$} \Comment{Builds
$\mtbody{f}, \mtbody{f'}$}
  \For {each $\{\pair{g,g'}\in \mapf \mid g,g'$ are reachable from $f,f'$\}}\label{step:forcheckCallequiv}
    \If {$\pair{g,g'} \in S$ or $\pair{g,g'}$ is marked \mtlabel} \label{step:substitute}
      \State \label{step:subst_uf}Replace calls to $g(\vec{in})$, $g'(\vec{in'})$, resp., with calls to \pneuf{g}, \pneufp{g'};
    \Else { inline $g,g'$ in their callers; } \label{step:isinline}
    \EndIf
  \EndFor
  \State \Return $\pair{f,f'}$;
\EndFunction
\end{algorithmic}
\end{minipage}
\caption{Isolation of functions: from $f$ to $\mtbody{f}$.}
\label{alg:isolate}
\end{algorithm} }

\subsubsection{\alg{isolate}.}
The function \alg{isolate} receives as input a pair $\pair{f,f'} \in \mapf$ and
a set $S$ of paired functions which, by construction (see
line~\ref{step:select}) contains only pairs from the same MSCCs as $f,f'$,
i.e., if $f \in m$ and $f' \in m'$, then $(g,g') \in S$ implies that $g \in m$
and $g' \in m'$. As output, it generates $\mtbody{f}$ and $\mtbody{f'}$, or
rather a relaxation thereof as explained after Eq.~(\ref{eq:isolation}). We
will occasionally refer to them as \emph{side 0} and \emph{side 1}. These
functions do not have function calls (other than to uninterpreted functions),
but may include inlined (nonrecursive) callees that were not proven to be
mutually terminating.
\Long{\footnote{There is some redundancy in the listing of this
algorithm, as demonstrated by the following case: $f$ calls $g$, and $g$ calls
$h$ (assume that $g$ is the only function to call $h$), where $h$ is inlined in
line~\ref{step:isinline}, but then $g$ is replaced with an uninterpreted
function in line~\ref{step:subst_uf}, which makes the former inlining
redundant. Our implementation avoids such cases by reaching the same result by
only considering  functions for inlining if their callers have already been
dealt-with.}}


The implementations of \pneufname and \pneufpname appear in Fig.~\ref{fig:ufs1}, and are rather self-explanatory. Their main role is to check call-equivalence. This is done by checking that they are called with the same set of inputs. When $\pair{g,g'}$ is marked \equivlabel, \pneufname and \pneufpname emulate the \emph{same} uninterpreted function\Long{(i.e., to return the same output given the same input)}, i.e.,
\[\forall \vec{in}.\ \mbox{{\sc uf}}(g, \vec{in}) = \mbox{{\sc uf'}}(g', \vec{in})\;.\]
 %
 When $\pair{g,g'}$ is not marked \equivlabel, \pneufname and \pneufpname emulate two \emph{different} uninterpreted functions.



\begin{figure}
\begin{tabular}{l} \hline
\begin{minipage}{\linewidth}
\mbox{ }\newline
\begin{algorithmic}[1]
\Function{uf}{function index $g$, input parameters $\vec{in}$}\label{startREs}
\Comment{Called in side 0}
  \If {$\vec{in} \in params[g]$} \Return the output of the earlier call \pneuf{$g$};
  \EndIf
  \State $params[g] := params[g]\; \bigcup\; \vec{in}$; \label{step:record1}
  \State \Return a non-deterministic output;
  \EndFunction \newline
%
%
\Function{uf'}{function index $g'$, input parameters $\vec{in'}$}
\Comment{Called in side 1}
  \If {$\vec{in'} \in params[g']$} \Return the output of the earlier call \pneufp{$g'$};
  \EndIf
  \State $params[g'] := params[g']\; \bigcup\; \vec{in'}$;
  \If {$\vec{in'} \in params[g]$} \Comment{$\pairwtag{g} \in \mapf$}\label{step:check_g'_le_g}
     \If {$\pair{g,g'}$ is marked \equivlabel}
       \State \Return the output of the earlier call {\sc uf}(g, $\vec{in'}$);
     \EndIf
     \State \Return a non-deterministic output;
  \EndIf
  \State\label{step:force_failure2}{\bf assert(0)}; \Comment{Not call-equivalent: $params[g'] \not\subseteq params[g]$}
\EndFunction \label{endREs}
\end{algorithmic}
\end{minipage}
\mbox{ }  \\
\mbox{ }  \\
\hline
\end{tabular}
%
\caption{Functions {\sc uf} and {\sc uf'} emulate uninterpreted functions if instantiated with functions that are mapped to one another. They are part of the generated program $\delta$, as shown in \alg{CallEquiv} of Alg.~\ref{alg:ProveMT}. }
\label{fig:ufs1}
\end{figure}


\subsubsection{\CheckCallEquiv.}
Our implementation is based on the C model checker
\tool{CBMC}~\cite{CK03}, which enables us to fully automate the
check for call-equivalence. \tool{CBMC} is complete for bounded programs (i.e.,
loops and recursions are bounded), and, indeed, the program $\delta$ we build
in \CheckCallEquiv is of that nature. It simply calls $\mtbody{f}$,
$\mtbody{f'}$ (which, recall, have no loops or function calls by construction),
with the same nondeterministic value, and asserts in the end that the set of
calls in $f$ is included in the set of calls in $f'$ (the other direction is
checked in lines~\ref{step:check_g'_le_g},~\ref{step:force_failure2} of {\sc uf'}). Examples of such generated programs
that we checked with CBMC are available online in~\cite{EKSurl}.


\subsection{An example}
The following example demonstrates Alg.~\ref{alg:ProveMT}. Consider the call graphs in Fig.~\ref{fig:callgraph}.
%
\begin{figure}[here]
\begin{center}
\scalebox{0.8}{
\begin{minipage}{0.5\linewidth}
\begin{tikzpicture}
\tikzset{node distance = 1.5cm}
\Vertex[Math=true]{f_1}
\SOWE[Math=true](f_1){f_2}
\SOEA[Math=true, style={fill = black!30}](f_1){f_4}
\WE[Math=true](f_2){f_3}
\WE[Math=true, style={fill = black!30}](f_3){f_5}
\GraphLoop[dist=1.0cm,dir=NO,style={->, thick}](f_5)
\Edge[style={->}](f_1)(f_2)
\Edge[style={->}](f_1)(f_4)
\Edge[style={->}](f_2)(f_3)
\Edge[style={->, bend right}](f_2)(f_4)
\Edge[style={->, bend right}](f_4)(f_2)
\Edge[style={->}](f_3)(f_5)
\end{tikzpicture}
\end{minipage}
\hspace{1 cm}
\begin{minipage}{0.5\linewidth}
\begin{tikzpicture}
\tikzset{node distance = 1.5cm}
\Vertex[Math=true]{f'_1}
\SOWE[Math=true](f'_1){f'_2}
\SOEA[Math=true, style={fill = black!30}](f'_1){f'_4}
\WE[Math=true](f'_2){f'_3}
\WE[Math=true, style={fill = black!30}](f'_3){f'_5}
\SO[Math=true](f'_1){f'_6}
\GraphLoop[dist=1.0cm,dir=NO,style={->, thick}](f'_5)
\Edge[style={->}](f'_1)(f'_2)
\Edge[style={->}](f'_1)(f'_4)
\Edge[style={->}](f'_2)(f'_3)
\Edge[style={->, bend right}](f'_2)(f'_6)
\Edge[style={->, bend right}](f'_6)(f'_4)
\Edge[style={->, bend right}](f'_4)(f'_2)
\Edge[style={->}](f'_3)(f'_5)
\end{tikzpicture}
\end{minipage}
}
\end{center}
\caption{Call graphs of the input programs $P,P'$. Partially equivalent functions are gray.}
\label{fig:callgraph}
\end{figure}



Assume that \map{f_i,f'_i} for $i = 1,\ldots,5$, and that the functions
represented by gray nodes are known to be partially equivalent to their counterparts.
%
Line~\ref{step:generate} generates the following nodes of the MSCC DAGs\Long{(listed bottom-up, left-to-right)}: $MD=\{\{f_5\}$, $\{f_3\}$,$\{f_2,f_4\},\{f_1\}\}$; $MD'=\{\{f'_5\}$, $\{f'_3\}$,$\{f'_2,f'_4,f'_6\}$,$\{f'_1\}\}$.
\Long{
\begin{itemize}
	\item $MD=\{\{f_5\}$, $\{f_3\}$,$\{f_2,f_4\},\{f_1\}\}$;
	\item $MD'=\{\{f'_5\}$, $\{f'_3\}$,$\{f'_2,f'_4,f'_6\}$,$\{f'_1\}\}$.
\end{itemize}
}
The MSCC mapping $\mapm$ in line~\ref{step:generate_map} is naturally derived from
$\mapf$.

\begin{table}[h]%[bottom]
\begin{tabular}{|p{1.5 cm}|l|p{8.6cm}|c|} \hline
MSCCs & Pair & Description & Res.\\\hline
%
$\{f_5\}$, $\{f_5'\}$ & $\pair{f_5, f'_5}$ & In line~\ref{step:select} the only possible $S$ is $\pair{f_5,f_5'}$. \alg{isolate} replaces the recursive call to $f_5, f_5'$ with \pneufname, \pneufpname, respectively \emulatesame.
%The latter emulate the same uninterpreted functions since $\pair{f_5,f_5'}$ is marked \equivlabel.
Assume \CheckCallEquiv  returns \true. $\pair{f_5, f'_5}$ is marked \mtlabel  in line~\ref{step:markNonTrivial}. & \checkmark \\ \hline
%
$\{f_3\}$, $\{f_3'\}$ & $\pair{f_3, f'_3}$ & This is a case of trivial MSCCs, which is handled in lines~\ref{step:m1}--\ref{step:markTrivial}. \alg{isolate} replaces the calls to $f_5,f_5'$ with \pneufname, \pneufpname, respectively \emulatesame.
%, emulating the same functions again.
Assume \CheckCallEquiv  returns \false. & \ding{55} \\ \hline
%
$\{f_2,f_4\}$, & & In line~\ref{step:select} let $S = \{\pair{f_2, f_2'}, \pair{f_4, f_4'}\}$. & \\ \cline{2-4}
%
$\{f_2', f_4', f_6'\}$& $\pair{f_2,f_2'}$ & In $f_2$ calls to $f_3$ are inlined, and calls to $f_4, f_5$ are replaced with calls to \pneufname. In $f_2'$ calls to $f_3', f_6'$ are inlined, and calls to $f_4', f_5'$ are replaced with calls to \pneufpname\emulatesame.
%Each replacing function pair emulates a pair of the same functions.
Assume \CheckCallEquiv  returns \true. & $\checkmark^c$ \\  \cline{2-4}
%
& $\pair{f_4,f_4'}$ & In $f_4$, $f_4'$ calls to $f_2$, $f_2'$ are respectively replaced with calls to \pneufname, \pneufpname\emulatedifferent.
%, which emulate two different uninterpreted functions because $\pair{f_2, f'_2}$ is not labeled \equivlabel.
Assume \CheckCallEquiv returns \true. Now $\pair{f_2, f'_2}$ and $\pair{f_4, f'_4}$ are marked \mtlabel in line~\ref{step:forall2}. & $\checkmark$ \\ \hline
%
$\{f_1\}$, $\{f_1'\}$ & $\pair{f_1, f'_1 }$ & Again, a case of a trivial MSCC. Calls to $f_2$, $f_2'$ are respectively replaced with \pneufname, \pneufpname \emulatedifferent
%, which emulate two different functions
, while calls to $f_4$, $f_4'$ are replaced with \pneufname, \pneufpname, respectively \emulatesame.
%, which emulate the same uninterpreted functions.
Assume \CheckCallEquiv returns \true. $\pair{\{f_1\}, \{f'_1\}}$ is marked \mtlabel. & $\checkmark$ \\ \hline
\end{tabular}
\caption{Applying Alg.~\ref{alg:ProveMT} to the call graphs in
Fig.~\ref{fig:callgraph}. `\checkmark' means that the pair is marked \mtlabel,
`$\checkmark^c$' that it is marked conditionally (it becomes unconditional once
all other pairs in $S$ are marked as well), and `\ding{55}' that it is not
marked.
\emulatesame and \emulatedifferent denote that \pneufname, \pneufpname emulate the same, or, respectively, different, uninterpreted functions.} \label{tab:final_iterations}
\end{table}

The progress of the algorithm is listed in Table~\ref{tab:final_iterations}. The output in this case, based on assumptions about the results of the checks for call-equivalence that are mentioned in the table, is that the following pairs of functions are marked as \mtlabel: $\pair{f_5,f_5'}$, $\pair{f_2,f_2'}$, $\pair{f_4,f_4'}$, and $\pair{f_1,f_1'}$.
%
\subsection{Choosing a vertex feedback set deterministically} \label{sec:determinization}
In line~\ref{step:choose} the choice of the set $S$ is nondeterministic. Our
implementation determinizes it by solving a series of optimization problems. In
the worst case this amounts to trying all sets, which is exponential in the
size of the MSCC. Observe, however, that large MSCCs are rare in real programs
%(even MSCCs of size 5 are rarely seen),
and, indeed, this has never posed a
computational problem in our experiments.

Our objective is to find a maximal
set $S$ of function pairs, because the larger the set is, the more functions
are declared to be mutually terminating in case of success. Further, larger
sets imply fewer functions to inline, and hence the burden on \CheckCallEquiv
is expected to be smaller.
%
Our implementation solves this optimization problem via a reduction to a pseudo-Boolean formula, which is then solved by \tool{minisat+}~\cite{ES06}. Each function node $g$ in $m$ (and $m'$) is associated with a Boolean variable $v_g$, indicating whether it is part of $S$. The objective is thus to maximize the sum of these variables that are mapped (those that are unmapped cannot be in $S$ anyway). In addition, there is a variable $e_{ij}$ for each edge $(i,j)$, which is set to true iff neither $i$ nor $j$ is in $S$. By enforcing a transitive closure, we guarantee that if there is a cycle of edges set to true (i.e., a cycle in which none of the nodes is in $S$), then the self edges (e.g., $e_{i,i}$) are set to \true as well. We then prevent such cycles by setting them to \false.
Let \emph{mapped($m$)} denote the set of functions in $m$ that are mapped. The problem formulation appears in Fig.~\ref{fig:pbs}, and is rather self-explanatory.
In case the chosen set $S$ fails (i.e., one of the pairs in $S$ cannot be proven to be mutually terminating), we add its negation (see constraint \#6) and repeat.

\begin{figure}
maximize $S$: \begin{minipage}{5 cm} \[ \max\sum_{g \in \mbox{\emph{mapped}(m)}}  v_g \] \end{minipage}

\vspace{0.1 cm}
subject to the following constraints, for $M \in \{m,m'\}$:
%
\[
\begin{array}{p{5 cm}l}
\mbox{1. Unmapped nodes are not in $S$:} &  \forall g \in (M \setminus mapped(M)).\ \lnot v_g \\
\mbox{2. Defining the edges:} & \forall \{i, j \mid (i,j) \mbox { is an edge in $M$}\}.\ \lnot v_i \land \lnot v_j \rightarrow \ e_{ij} \\
\mbox{3. Transitive closure:} &  \forall 0 < i,j,k \leq |M|.\ e_{ij} \land e_{jk} \rightarrow e_{ik} \\
\mbox{4. Self loops are not allowed:} & \forall 0 < i \leq |M|.\ \lnot e_{ii} \\
\mbox{5. Enforce mapping:} & \forall \pair{g,g'} \in \mapf, g \in m.\ v_g \leftrightarrow v_g' \\
\mbox{6. For each  failed solution $Sl$:} & \bigvee_{\pair{g,g'} \in Sl} \lnot v_g \\
\end{array}
\]
%
\caption{A pseudo-Boolean formulation of the optimization problem of finding the largest set of function pairs from $m,m'$ that intersect all cycles in both $m$ and $m'$.}\label{fig:pbs}
\end{figure}
%

\Long{
Three optimizations to this process that we implemented in \tool{rvt} are:
\begin{itemize}
\item \emph{Recycling proofs}: When failing with a set $S$, we save the strict subset $s \subset S$ for which we were able to prove call-equivalence. Let $S'$ be a new set under consideration and let $g \in S' \cap s$. Denote by $\mtbody{g}_S$ the function $\mtbody{g}$ as constructed given the set $S$. Then the positive result of $g$ can be reused if $\mtbody{g}_S$ is equally or more abstract than $\mtbody{g}_{S'}$. This condition holds if none of the functions inlined in $\mtbody{g}_S$ are abstracted in $\mtbody{g}_{S'}$.

\item \emph{Generalizing counterexamples}: When failing with a set $S$, it is frequently possible to  find a strict subset of $S$ that is sufficient to make the proof fail. Specifically, the subset of $S$ equal to the failing node plus nodes in $S$ that can be reached from it not through any other node in $S$, constitute such a subset. As an example, consider the MSCC $\{f \leftrightarrow g, g \leftrightarrow h \}$, and $S = \{f,g,h\}$. If $f$ fails, then we must remove from $S$ either $f$ or	 $g$, regardless of what we do with $h$. Hence here we can return the subset $\{f,g\}$ as the failing set, rather than $S$ itself, which strengthens constraint \#6.
    %In other words, adding the negation of this subset removes from the solution space sets that are bound to fail.

\item \emph{Higher priority for \pe pairs} Among equal-sized sets, it is better to give priority to pairs that are marked as \equivlabel, because there is a better chance of success with such functions (recall that these are replaced with the same uniniterpreted functions). This strategy can be implemented by giving higher weight to such functions in the objective.
\end{itemize}
}

%\section{A decomposition algorithm}
\label{sec:decomposition}

In this section we present an algorithm for proving mutual termination of full programs.
As mentioned in Sec.~\ref{sec:rule}, the call graph of a program can be viewed
as a DAG where the nodes correspond to MSCCs. After building a mapping between the MSCCs of the two call graphs, the algorithm traverses the DAG
bottom-up. For each mapped pair of MSCCs $m,m'$, it attempts to prove the mutual termination of their mapped functions, based on \mtermdp.

The algorithm is inspired by a similar algorithm for verification of \emph{partial
equivalence}, which is described in a technical report~\cite{GS11}
\Long{\footnote{\emph{Note to reviewers}: this report is also currently under review
for a journal~\cite{GS10}}}. The algorithm here is more involved, however,
because it handles differently cases in which the checked functions are also
partially equivalent (recall that this information, i.e., which functions are
known to be partially equivalent, is part of the input to the algorithm).
Furthermore, the algorithm in~\cite{GS11} is described with a non-deterministic
step, and here we suggest a method for determinizing it.

The preprocessing and mapping is as in Sec.~\ref{sec:base}. Hence the program is loop-free, globals accessed by a function are sent instead as additional inputs, and there is a (possibly partial) mapping $\mapf$ between the functions of $P$ and $P'$.

%Note that wrong mapping does not affect soundness: it is used for generating the verification conditions. Hence, wrong mapping can only cause a proof failure.

%For a simpler description, the algorithm assumes that each function receives only one parameter. But this is not a limitation: we can consider a function that receives several parameters as if it gets a tuple, i.e., one complex parameter. Thus, the aforementioned expectation of the same input parameters amounts to an expectation of equivalent tuple structures among the members of each function pair.

\subsection{The algorithm}
\Long{
Table~\ref{tab:markings} contains the labels we use in the decomposition
algorithm (Alg.~\ref{alg:ProveMT}).

\begin{table}[here]
\begin{tabular}{||l|l|p{6.2cm}||} \hline \label{table:markings}
\textbf{Predicate label}	& \textbf{What labeled}  & \textbf{Meaning} \\ \hline
\mtlabel      & \pairtag{f}$ \in\mapf$ & The mapped functions \comandtag{f} are mutually terminating \\ \hline
\equivlabel   & \pairtag{f}$ \in\mapf$ & The mapped functions \comandtag{f} are partially equivalent \\ \hline
\mscccovered  & \pairtag{m}$ \in\mapm$ & The MSCC pair has been processed. \\ \hline
\end{tabular}
\caption{Predicate labels used in the decomposition algorithm.}
\label{tab:markings}
\end{table}
}
The input to Alg.~\ref{alg:ProveMT} is \comtag{P}, a (possibly partial) mapping $\mapf$ between their functions, and (implicitly) those paired functions that are known to be partially equivalent. Its output is a set of function pairs that are marked as \mtlabel, indicating it succeeded to prove their mutual termination based on \mtermdp. We now describe the three functions used by this algorithm.


\subsubsection{\ProveMT.} \label{sec:ProveMT}
This entry function traverses the call
graphs of $P,P'$ bottom-up, each time focusing on a pair of MSCCs. In
line~\ref{step:inline} it inlines all nonrecursive functions that are not
mapped. In line $\ref{step:rename}$ it uses renaming to resolve possible name
collisions between the globals of the two input programs. The next line builds
the MSCC DAGs $MD$ and $MD'$ from the call graphs, as explained in
Sec.\ref{sec:rule}. Line~\ref{step:generate_map} attempts to build a bijective
mapping $\mapm$ between the nodes of $MD$ and $MD'$ which is consistent with
$map_f$. Namely, if $\pair{m,m'} \in \mapm$, $f$ is a function in $m$, and
$\pair{f,f'} \in \mapf$, then $f'$ is a function in $m'$ (and vice-versa). If
such a mapping is impossible, the algorithm aborts. In practice one may run the
algorithm bottom-up until reaching nonmapped MSCCs, but we omit this option
here for brevity.

The bottom-up traversal starts in line~\ref{step:while}. Initially all MSCCs
are unmarked. The algorithm searches for a next unmarked pair \pairtag{m} of
MSCCs such that all its children pairs are marked. If \comtag{m} are trivial
(see Sec.~\ref{sec:rule} for a definition), then line~\ref{step:Check} simply
checks the call-equivalence of the function pair \pairtag{f} that constitutes
\pairtag{m}, and marks them accordingly in line~\ref{step:markTrivial}. 
Note that even if the descendants of \comtag{m} are mutually-terminating, \comtag{m}
are not necessarily so, because they may call their descendants with different
parameters. Also note that if this check fails, we continue to check their ancestors, (in contrast to the case of non-trivial MSCCs listed next), because even if \pairtag{f} are not
mutually terminating for every input, their callers may still be (they can be mutually
terminating in the context of their callers). We can check this by inlining
them, which is only possible because they are not recursive.



\noindent
\begin{algorithm}
\begin{minipage}{\linewidth}
\begin{algorithmic}[1]
\Function{ProveMT}{Programs \comtag{P}, map between functions $\mapf$}
\State \label{step:inline} Inline non-recursive non-mapped functions;
\State \label{step:rename} Solve name collisions in global identifiers of $P,P'$ by
renaming.
\State \label{step:generate} Generate MSCC DAGs \comtag{MD} from the
call graphs of $P,P'$;

\State If possible, generate a bijective map
\label{step:generate_map} $\mapm$ between the nodes of $MD$ \newline \mbox{~~~~~~} and $MD'$ that is consistent with $\mapf$;
\Long{\footnote{It is desirable but not necessary to add pairs of trivial nodes to $\mapm$.}}
Otherwise \abort.

\While {$\exists \pair{m,m'} \in \mapm$
          not marked \mscccovered but its children are,} \label{step:while}
   \State \label{step:choose} Choose such a pair $\pair{m,m'} \in \mapm$ and mark it \mscccovered
   \If {$m,m'$ are trivial} \label{step:m1}
     \State Let $f,f'$ be the functions in $m,m'$, respectively;
     \If {\CheckCallEquiv(\alg{isolate}$(f,f',\emptyset)$)} \label{step:Check}
     mark $f,f'$ as \mtlabel; \label{step:markTrivial}
     \EndIf
   \Else
    \label{step:select} Select non-deterministically
         $S \subseteq \{\pair{f,f'}\mid \pair{f,f'}\in\mapf(m)\}$ \newline
          \mbox{~~~~~~~~~~~~~~~~} that intersect every cycle in \comandtag{m};
      \For {each $\pair{f,f'} \in S$}\label{step:forall}
        \If  {$\lnot$\CheckCallEquiv(\alg{isolate}$(f,f',S)$)} \label{step:abort}{mark ancestors as \mscccovered;}
        \EndIf
      \EndFor
      \For {each $\pair{f,f'} \in S$} \label{step:forall2}
         mark $f,f'$ as \mtlabel; \label{step:markNonTrivial}
      \EndFor
   \EndIf
\EndWhile
\EndFunction%\newline

\vspace{0.25 cm}


\Function{isolate}{functions \comtag{f}, function pairs $S$} \Comment{Builds
$\mtbody{f}, \mtbody{f'}$}
  \For {each $\{\pair{g,g'}\in \mapf \mid g,g'$ are reachable from $f,f'$\}}\label{step:forcheckCallequiv}
    \If {$\pair{g,g'} \in S$ or $\pair{g,g'}$ is marked \mtlabel} \label{step:substitute}
      \State \label{step:subst_uf}Replace calls to $g(\vec{in})$, $g'(\vec{in'})$ with calls to \pneuf{g}, \pneufp{g'}, resp.;
    \Else { inline $g,g'$ in their callers; } \label{step:isinline}
    \EndIf
  \EndFor
  \State \Return $\pair{f,f'}$;
\EndFunction
\vspace{0.25 cm}

\Function {CallEquiv}{A pair of isolated functions $\pair{\mtbody{f}, \mtbody{f'}}$}
\label{step:checkCallEquiv}
  \State \label{step:gen_main} Let $\delta$ denote the program:
  \[
  \begin{array}{l}
 \rhd \mbox{ here add the definitions of {\sc uf()} and {\sc uf'()} (see Fig.~\ref{fig:ufs1}).}\\
 \vec{in} := nondet(); \mtbody{f}(\vec{in}); \mtbody{f'}(\vec{in}); \\
% \mbox{assert}(params(${\sc UF}$) \subseteq params(${\sc UF'}$));\\
 \forall \pair{g,g'} \mbox{ called in } f,f'.\ \mbox{assert}(params[g] \subseteq params[g']);\\
 \end{array}
 \]
  \State \Return CBMC($\delta$); \label{step:cbmc}
\EndFunction

\end{algorithmic}
\end{minipage}
\caption{Pseudo-code for a bottom-up decomposition algorithm for proving that pairs of functions mutually terminate.}
\label{alg:ProveMT}
\end{algorithm}


Next, consider the case that the selected \comtag{m} in line~\ref{step:choose}
are not trivial. In line~\ref{step:select} the algorithm chooses
non-deterministically a subset $S$ of pairs from $\mapf(m)$ that intersects all
the cycles in \comandtag{m}\Long{\footnote{In graph-theoretic terms, the functions in
$S$ constitute a \emph{feedback vertex set} of both $m$ and $m'$\Long{~\cite{E79}}.}}.
This guarantees that we can always inline the functions in \comtag{m}
that are not in $S$. Determinization of this step will be considered in
subsection~\ref{sec:determinization}. If \CheckCallEquiv returns \true for all
the function pairs in $S$, then all those pairs are labeled as \mtlabel in
line~\ref{step:forall2}. Otherwise it abandons the attempt to prove their ancestors in line~\ref{step:abort}: it cannot prove
that mapped functions in $\pair{m,m'}$ are mutually
terminating, nor can it inline these functions in their callers, so we cannot check all its ancestors.

Regardless of whether \pairtag{m} are trivial, they get marked as \coverlabel
in line~\ref{step:choose}, and the loop in \ProveMT  continues to another
pair.




\comment{
\begin{algorithm}[t]
\begin{minipage}{\linewidth}
\begin{algorithmic}[1]

\Function{isolate}{functions \comtag{f}, function pairs $S$} \Comment{Builds
$\mtbody{f}, \mtbody{f'}$}
  \For {each $\{\pair{g,g'}\in \mapf \mid g,g'$ are reachable from $f,f'$\}}\label{step:forcheckCallequiv}
    \If {$\pair{g,g'} \in S$ or $\pair{g,g'}$ is marked \mtlabel} \label{step:substitute}
      \State \label{step:subst_uf}Replace calls to $g(\vec{in})$, $g'(\vec{in'})$, resp., with calls to \pneuf{g}, \pneufp{g'};
    \Else { inline $g,g'$ in their callers; } \label{step:isinline}
    \EndIf
  \EndFor
  \State \Return $\pair{f,f'}$;
\EndFunction
\end{algorithmic}
\end{minipage}
\caption{Isolation of functions: from $f$ to $\mtbody{f}$.}
\label{alg:isolate}
\end{algorithm} }

\subsubsection{\alg{isolate}.}
The function \alg{isolate} receives as input a pair $\pair{f,f'} \in \mapf$ and
a set $S$ of paired functions which, by construction (see
line~\ref{step:select}) contains only pairs from the same MSCCs as $f,f'$,
i.e., if $f \in m$ and $f' \in m'$, then $(g,g') \in S$ implies that $g \in m$
and $g' \in m'$. As output, it generates $\mtbody{f}$ and $\mtbody{f'}$, or
rather a relaxation thereof as explained after Eq.~(\ref{eq:isolation}). We
will occasionally refer to them as \emph{side 0} and \emph{side 1}. These
functions do not have function calls (other than to uninterpreted functions),
but may include inlined (nonrecursive) callees that were not proven to be
mutually terminating.
\Long{\footnote{There is some redundancy in the listing of this
algorithm, as demonstrated by the following case: $f$ calls $g$, and $g$ calls
$h$ (assume that $g$ is the only function to call $h$), where $h$ is inlined in
line~\ref{step:isinline}, but then $g$ is replaced with an uninterpreted
function in line~\ref{step:subst_uf}, which makes the former inlining
redundant. Our implementation avoids such cases by reaching the same result by
only considering  functions for inlining if their callers have already been
dealt-with.}}


The implementations of \pneufname and \pneufpname appear in Fig.~\ref{fig:ufs1}, and are rather self-explanatory. Their main role is to check call-equivalence. This is done by checking that they are called with the same set of inputs. When $\pair{g,g'}$ is marked \equivlabel, \pneufname and \pneufpname emulate the \emph{same} uninterpreted function\Long{(i.e., to return the same output given the same input)}, i.e.,
\[\forall \vec{in}.\ \mbox{{\sc uf}}(g, \vec{in}) = \mbox{{\sc uf'}}(g', \vec{in})\;.\]
 %
 When $\pair{g,g'}$ is not marked \equivlabel, \pneufname and \pneufpname emulate two \emph{different} uninterpreted functions.



\begin{figure}
\begin{tabular}{l} \hline
\begin{minipage}{\linewidth}
\mbox{ }\newline
\begin{algorithmic}[1]
\Function{uf}{function index $g$, input parameters $\vec{in}$}\label{startREs}
\Comment{Called in side 0}
  \If {$\vec{in} \in params[g]$} \Return the output of the earlier call \pneuf{$g$};
  \EndIf
  \State $params[g] := params[g]\; \bigcup\; \vec{in}$; \label{step:record1}
  \State \Return a non-deterministic output;
  \EndFunction \newline
%
%
\Function{uf'}{function index $g'$, input parameters $\vec{in'}$}
\Comment{Called in side 1}
  \If {$\vec{in'} \in params[g']$} \Return the output of the earlier call \pneufp{$g'$};
  \EndIf
  \State $params[g'] := params[g']\; \bigcup\; \vec{in'}$;
  \If {$\vec{in'} \in params[g]$} \Comment{$\pairwtag{g} \in \mapf$}\label{step:check_g'_le_g}
     \If {$\pair{g,g'}$ is marked \equivlabel}
       \State \Return the output of the earlier call {\sc uf}(g, $\vec{in'}$);
     \EndIf
     \State \Return a non-deterministic output;
  \EndIf
  \State\label{step:force_failure2}{\bf assert(0)}; \Comment{Not call-equivalent: $params[g'] \not\subseteq params[g]$}
\EndFunction \label{endREs}
\end{algorithmic}
\end{minipage}
\mbox{ }  \\
\mbox{ }  \\
\hline
\end{tabular}
%
\caption{Functions {\sc uf} and {\sc uf'} emulate uninterpreted functions if instantiated with functions that are mapped to one another. They are part of the generated program $\delta$, as shown in \alg{CallEquiv} of Alg.~\ref{alg:ProveMT}. }
\label{fig:ufs1}
\end{figure}


\subsubsection{\CheckCallEquiv.}
Our implementation is based on the C model checker
\tool{CBMC}~\cite{CK03}, which enables us to fully automate the
check for call-equivalence. \tool{CBMC} is complete for bounded programs (i.e.,
loops and recursions are bounded), and, indeed, the program $\delta$ we build
in \CheckCallEquiv is of that nature. It simply calls $\mtbody{f}$,
$\mtbody{f'}$ (which, recall, have no loops or function calls by construction),
with the same nondeterministic value, and asserts in the end that the set of
calls in $f$ is included in the set of calls in $f'$ (the other direction is
checked in lines~\ref{step:check_g'_le_g},~\ref{step:force_failure2} of {\sc uf'}). Examples of such generated programs
that we checked with CBMC are available online in~\cite{EKSurl}.


\subsection{An example}
The following example demonstrates Alg.~\ref{alg:ProveMT}. Consider the call graphs in Fig.~\ref{fig:callgraph}.
%
\begin{figure}[here]
\begin{center}
\scalebox{0.8}{
\begin{minipage}{0.5\linewidth}
\begin{tikzpicture}
\tikzset{node distance = 1.5cm}
\Vertex[Math=true]{f_1}
\SOWE[Math=true](f_1){f_2}
\SOEA[Math=true, style={fill = black!30}](f_1){f_4}
\WE[Math=true](f_2){f_3}
\WE[Math=true, style={fill = black!30}](f_3){f_5}
\GraphLoop[dist=1.0cm,dir=NO,style={->, thick}](f_5)
\Edge[style={->}](f_1)(f_2)
\Edge[style={->}](f_1)(f_4)
\Edge[style={->}](f_2)(f_3)
\Edge[style={->, bend right}](f_2)(f_4)
\Edge[style={->, bend right}](f_4)(f_2)
\Edge[style={->}](f_3)(f_5)
\end{tikzpicture}
\end{minipage}
\hspace{1 cm}
\begin{minipage}{0.5\linewidth}
\begin{tikzpicture}
\tikzset{node distance = 1.5cm}
\Vertex[Math=true]{f'_1}
\SOWE[Math=true](f'_1){f'_2}
\SOEA[Math=true, style={fill = black!30}](f'_1){f'_4}
\WE[Math=true](f'_2){f'_3}
\WE[Math=true, style={fill = black!30}](f'_3){f'_5}
\SO[Math=true](f'_1){f'_6}
\GraphLoop[dist=1.0cm,dir=NO,style={->, thick}](f'_5)
\Edge[style={->}](f'_1)(f'_2)
\Edge[style={->}](f'_1)(f'_4)
\Edge[style={->}](f'_2)(f'_3)
\Edge[style={->, bend right}](f'_2)(f'_6)
\Edge[style={->, bend right}](f'_6)(f'_4)
\Edge[style={->, bend right}](f'_4)(f'_2)
\Edge[style={->}](f'_3)(f'_5)
\end{tikzpicture}
\end{minipage}
}
\end{center}
\caption{Call graphs of the input programs $P,P'$. Partially equivalent functions are gray.}
\label{fig:callgraph}
\end{figure}



Assume that \map{f_i,f'_i} for $i = 1,\ldots,5$, and that the functions
represented by gray nodes are known to be partially equivalent to their counterparts.
%
Line~\ref{step:generate} generates the following nodes of the MSCC DAGs\Long{(listed bottom-up, left-to-right)}: $MD=\{\{f_5\}$, $\{f_3\}$,$\{f_2,f_4\},\{f_1\}\}$; $MD'=\{\{f'_5\}$, $\{f'_3\}$,$\{f'_2,f'_4,f'_6\}$,$\{f'_1\}\}$.
\Long{
\begin{itemize}
	\item $MD=\{\{f_5\}$, $\{f_3\}$,$\{f_2,f_4\},\{f_1\}\}$;
	\item $MD'=\{\{f'_5\}$, $\{f'_3\}$,$\{f'_2,f'_4,f'_6\}$,$\{f'_1\}\}$.
\end{itemize}
}
The MSCC mapping $\mapm$ in line~\ref{step:generate_map} is naturally derived from
$\mapf$.

\begin{table}[h]%[bottom]
\begin{tabular}{|p{1.5 cm}|l|p{8.6cm}|c|} \hline
MSCCs & Pair & Description & Res.\\\hline
%
$\{f_5\}$, $\{f_5'\}$ & $\pair{f_5, f'_5}$ & In line~\ref{step:select} the only possible $S$ is $\pair{f_5,f_5'}$. \alg{isolate} replaces the recursive call to $f_5, f_5'$ with \pneufname, \pneufpname, respectively \emulatesame.
%The latter emulate the same uninterpreted functions since $\pair{f_5,f_5'}$ is marked \equivlabel.
Assume \CheckCallEquiv  returns \true. $\pair{f_5, f'_5}$ is marked \mtlabel  in line~\ref{step:markNonTrivial}. & \checkmark \\ \hline
%
$\{f_3\}$, $\{f_3'\}$ & $\pair{f_3, f'_3}$ & This is a case of trivial MSCCs, which is handled in lines~\ref{step:m1}--\ref{step:markTrivial}. \alg{isolate} replaces the calls to $f_5,f_5'$ with \pneufname, \pneufpname, respectively \emulatesame.
%, emulating the same functions again.
Assume \CheckCallEquiv  returns \false. & \ding{55} \\ \hline
%
$\{f_2,f_4\}$, & & In line~\ref{step:select} let $S = \{\pair{f_2, f_2'}, \pair{f_4, f_4'}\}$. & \\ \cline{2-4}
%
$\{f_2', f_4', f_6'\}$& $\pair{f_2,f_2'}$ & In $f_2$ calls to $f_3$ are inlined, and calls to $f_4, f_5$ are replaced with calls to \pneufname. In $f_2'$ calls to $f_3', f_6'$ are inlined, and calls to $f_4', f_5'$ are replaced with calls to \pneufpname\emulatesame.
%Each replacing function pair emulates a pair of the same functions.
Assume \CheckCallEquiv  returns \true. & $\checkmark^c$ \\  \cline{2-4}
%
& $\pair{f_4,f_4'}$ & In $f_4$, $f_4'$ calls to $f_2$, $f_2'$ are respectively replaced with calls to \pneufname, \pneufpname\emulatedifferent.
%, which emulate two different uninterpreted functions because $\pair{f_2, f'_2}$ is not labeled \equivlabel.
Assume \CheckCallEquiv returns \true. Now $\pair{f_2, f'_2}$ and $\pair{f_4, f'_4}$ are marked \mtlabel in line~\ref{step:forall2}. & $\checkmark$ \\ \hline
%
$\{f_1\}$, $\{f_1'\}$ & $\pair{f_1, f'_1 }$ & Again, a case of a trivial MSCC. Calls to $f_2$, $f_2'$ are respectively replaced with \pneufname, \pneufpname \emulatedifferent
%, which emulate two different functions
, while calls to $f_4$, $f_4'$ are replaced with \pneufname, \pneufpname, respectively \emulatesame.
%, which emulate the same uninterpreted functions.
Assume \CheckCallEquiv returns \true. $\pair{\{f_1\}, \{f'_1\}}$ is marked \mtlabel. & $\checkmark$ \\ \hline
\end{tabular}
\caption{Applying Alg.~\ref{alg:ProveMT} to the call graphs in
Fig.~\ref{fig:callgraph}. `\checkmark' means that the pair is marked \mtlabel,
`$\checkmark^c$' that it is marked conditionally (it becomes unconditional once
all other pairs in $S$ are marked as well), and `\ding{55}' that it is not
marked.
\emulatesame and \emulatedifferent denote that \pneufname, \pneufpname emulate the same, or, respectively, different, uninterpreted functions.} \label{tab:final_iterations}
\end{table}

The progress of the algorithm is listed in Table~\ref{tab:final_iterations}. The output in this case, based on assumptions about the results of the checks for call-equivalence that are mentioned in the table, is that the following pairs of functions are marked as \mtlabel: $\pair{f_5,f_5'}$, $\pair{f_2,f_2'}$, $\pair{f_4,f_4'}$, and $\pair{f_1,f_1'}$.
%
\subsection{Choosing a vertex feedback set deterministically} \label{sec:determinization}
In line~\ref{step:choose} the choice of the set $S$ is nondeterministic. Our
implementation determinizes it by solving a series of optimization problems. In
the worst case this amounts to trying all sets, which is exponential in the
size of the MSCC. Observe, however, that large MSCCs are rare in real programs
%(even MSCCs of size 5 are rarely seen),
and, indeed, this has never posed a
computational problem in our experiments.

Our objective is to find a maximal
set $S$ of function pairs, because the larger the set is, the more functions
are declared to be mutually terminating in case of success. Further, larger
sets imply fewer functions to inline, and hence the burden on \CheckCallEquiv
is expected to be smaller.
%
Our implementation solves this optimization problem via a reduction to a pseudo-Boolean formula, which is then solved by \tool{minisat+}~\cite{ES06}. Each function node $g$ in $m$ (and $m'$) is associated with a Boolean variable $v_g$, indicating whether it is part of $S$. The objective is thus to maximize the sum of these variables that are mapped (those that are unmapped cannot be in $S$ anyway). In addition, there is a variable $e_{ij}$ for each edge $(i,j)$, which is set to true iff neither $i$ nor $j$ is in $S$. By enforcing a transitive closure, we guarantee that if there is a cycle of edges set to true (i.e., a cycle in which none of the nodes is in $S$), then the self edges (e.g., $e_{i,i}$) are set to \true as well. We then prevent such cycles by setting them to \false.
Let \emph{mapped($m$)} denote the set of functions in $m$ that are mapped. The problem formulation appears in Fig.~\ref{fig:pbs}, and is rather self-explanatory.
In case the chosen set $S$ fails (i.e., one of the pairs in $S$ cannot be proven to be mutually terminating), we add its negation (see constraint \#6) and repeat.

\begin{figure}
maximize $S$: \begin{minipage}{5 cm} \[ \max\sum_{g \in \mbox{\emph{mapped}(m)}}  v_g \] \end{minipage}

\vspace{0.1 cm}
subject to the following constraints, for $M \in \{m,m'\}$:
%
\[
\begin{array}{p{5 cm}l}
\mbox{1. Unmapped nodes are not in $S$:} &  \forall g \in (M \setminus mapped(M)).\ \lnot v_g \\
\mbox{2. Defining the edges:} & \forall \{i, j \mid (i,j) \mbox { is an edge in $M$}\}.\ \lnot v_i \land \lnot v_j \rightarrow \ e_{ij} \\
\mbox{3. Transitive closure:} &  \forall 0 < i,j,k \leq |M|.\ e_{ij} \land e_{jk} \rightarrow e_{ik} \\
\mbox{4. Self loops are not allowed:} & \forall 0 < i \leq |M|.\ \lnot e_{ii} \\
\mbox{5. Enforce mapping:} & \forall \pair{g,g'} \in \mapf, g \in m.\ v_g \leftrightarrow v_g' \\
\mbox{6. For each  failed solution $Sl$:} & \bigvee_{\pair{g,g'} \in Sl} \lnot v_g \\
\end{array}
\]
%
\caption{A pseudo-Boolean formulation of the optimization problem of finding the largest set of function pairs from $m,m'$ that intersect all cycles in both $m$ and $m'$.}\label{fig:pbs}
\end{figure}
%

\Long{
Three optimizations to this process that we implemented in \tool{rvt} are:
\begin{itemize}
\item \emph{Recycling proofs}: When failing with a set $S$, we save the strict subset $s \subset S$ for which we were able to prove call-equivalence. Let $S'$ be a new set under consideration and let $g \in S' \cap s$. Denote by $\mtbody{g}_S$ the function $\mtbody{g}$ as constructed given the set $S$. Then the positive result of $g$ can be reused if $\mtbody{g}_S$ is equally or more abstract than $\mtbody{g}_{S'}$. This condition holds if none of the functions inlined in $\mtbody{g}_S$ are abstracted in $\mtbody{g}_{S'}$.

\item \emph{Generalizing counterexamples}: When failing with a set $S$, it is frequently possible to  find a strict subset of $S$ that is sufficient to make the proof fail. Specifically, the subset of $S$ equal to the failing node plus nodes in $S$ that can be reached from it not through any other node in $S$, constitute such a subset. As an example, consider the MSCC $\{f \leftrightarrow g, g \leftrightarrow h \}$, and $S = \{f,g,h\}$. If $f$ fails, then we must remove from $S$ either $f$ or	 $g$, regardless of what we do with $h$. Hence here we can return the subset $\{f,g\}$ as the failing set, rather than $S$ itself, which strengthens constraint \#6.
    %In other words, adding the negation of this subset removes from the solution space sets that are bound to fail.

\item \emph{Higher priority for \pe pairs} Among equal-sized sets, it is better to give priority to pairs that are marked as \equivlabel, because there is a better chance of success with such functions (recall that these are replaced with the same uniniterpreted functions). This strategy can be implemented by giving higher weight to such functions in the objective.
\end{itemize}
}


\section{An inference rule for proving termination} \label{sec:term}
We now consider a different variant of the mutual termination problem:
%\begin{quote}
\emph{Given that a program $P$ terminates, does $P'$ terminate as well? }
%\end{quote}
%
Clearly this problem can be reduced to that of mutual termination, but in fact it can also
be solved with a weaker premise. We first define $term(f)$ to denote that $f$ terminates and
\[
\begin{array}{l}
  \ \ \ \ \\
\callc(f,f') \doteq \\
  \ \ \ \forall \text{\pairtag{g}} \in \mapf, in_f, in_g.\ g'(in_g) \in calls (f'(in_f)) \Rightarrow g(in_g) \in calls(f(in_f))\;.\\
  \ \ \ \
\end{array}
\]
%\begin{definition}[Call-containment of functions]
%Define \emph{call-containment} of $f$ in $f'$ as following
%$f$ \emph{call-contains} in $f'$  if and only if
%\[
%\forall \text{\pairtag{g}} \in \mapf, in_f, in_g.\ g'(in_g) \in calls (f'(in_f)) \Rightarrow g(in_g) \in calls(f(in_f))\;. \]
%\end{definition}
%Denote by $\callc(f,f')$ that $f$ call-contains $f'$.
\newline Using these predicates, we can now define the rule for leaf MSCCs \comtag{m}:
\begin{equation}
\frac
{\forall\left\langle f, f' \right\rangle\in\mapf(m).\left(term(f)\wedge \callc(\upbody{f}, \upbody{f'}\right)}
{\forall\left\langle f, f' \right\rangle\in\mapf(m).~term(f')}
\mbox{\ \termd} \;.
\end{equation}
%A soundness proof appears in Appendix~\ref{sec:proof2}.

\begin{theorem} \termd is sound. \end{theorem}
\begin{proof}
The proof follows similar lines to that of \mtermdp. We give a proof sketch.
Falsely assume that there is a function $f'$ in $m'$ that does not terminate, whereas for all $\pair{g,g'} \in \mapf(m)$, $\callc(g,g')$.
There exists a value $in$ such that $f'(in)$ does not terminate. The infinite call stack of $f'(in)$ must contain a call, say from $h'(in_1)$ to $g'(in_2)$, whereas $h(in_1)$ does not call $g(in_2)$ in the call stack of $f(in)$ (assuming $\{$\pairtag{g}, \pairtag{h}$\} \subseteq \mapf$).
This contradicts our premise that $\callc(h,h')$ is true. \qed
\end{proof}


%
Note that call-equivalence (Def.~\ref{def:re_func}) is simply bi-directional
call-containment.
%, which implies that the premise of \termd is weaker than that of \mtermd (see~(\ref{eq:new})) for proving mutual termination. Observe that a
A generalization to non-leaf MSCCs can be done as in (\ref{eq:new1}).

The decomposition algorithm of the previous section (Alg.~\ref{alg:ProveMT}) applies with the following
change: the last statement of line~\ref{step:gen_main} (asserting $params[g] \subseteq
params[g']$) should be removed. The only assertion that should be verified is thus inside
\pneufname (line \ref{step:force_failure2} in Fig.~\ref{fig:ufs1}), which checks that every call on side 0 is matched by a call on side 1.

\section{Experience and conclusions} \label{sec:experiments}
We implemented Alg.~\ref{alg:ProveMT} in RVT~\cite{GS09,EKSurl}, and tested it with many small programs and one real software project, as follows.

\subsubsection{Small programs.}
Our benchmark suite of small programs\footnote{The programs are available online in~\cite{EKSurl}} includes: (a) various small programs that check different cases handled by Alg.~\ref{alg:ProveMT}, (b) 11 programs from the \tool{rvt}
regression test suite, which on the one hand were designed
for testing various features of RVT, including non-isomorphic call graphs and
complex MSCCs, and on the other are interesting for checking mutual termination, (c) various famous termination problems, such as McCarthy91, Collatz, the Ackermann function and more (see Example %s~\ref{ex:mccarthy} and~
\ref{ex:collatz}), and (d) each of the (single-threaded) examples
from Cook et al.~\cite{CPR11}, where each program was checked against a variation thereof that we attained by
modifying the code in a non-trivial manner.

Since the only exponential component in Alg.~\ref{alg:ProveMT} is the call to
CBMC (and also choosing $S$ - but the MSCCs we encountered were not large, and
never took a significant amount of time), we do not list here the run times.
Run time has been a problem in very few cases, all of which due to line~\ref{step:cbmc} in Alg.~\ref{alg:ProveMT} -- the call to CBMC. Such performance problems occurred in our experiments either when the program included difficult operators (such as multiplication  over
32 bit integers) that RVT could not abstract, or when the compared functions
received recursive data structures as input. In such cases RVT sends a
nondeterministic yet equal data structure unrolled $k$ times, where $k$ is user-defined, as explained in~\cite{GS09}. This unrolling sometimes causes performance problems to CBMC.


\subsubsection{A real software project.}
We also tested our tool on the open source project \alg{Betik}~\cite{Betik}, which is an interpreter for a scripting language. The code has 2 -- 2.5 KLOC (depending on the version). It has many loops and recursive functions, including mutual recursion forming an MSCC of size 14. % (large MSCCs are typical in parsers).
%, as well as advanced C features such as function pointers, heap allocation [*?*], complex data structures and unions.
%
We compared eight consecutive versions of this program from the code repository, i.e., seven comparisons. The amount of changes between the versions varied with an average of 3--4 (related) functions. Somewhat to our surprise, many of the changes do \emph{not} preserve termination behavior in a free context, mostly because these functions traverse global data structures on the heap.

In five out of the seven comparisons, RVT discovered correctly, in less than 2 minutes each, that the programs contained mapped functions that do not mutually terminate.
An example is a function
\Long{Fig.~\ref{fig:int_value} displays two versions of a function} called \alg{int\_value()}, which receives a pointer to a node in a syntax tree.
The old version compared the type of the node to several values, and if none of them
matched it simply returned the input node. In the new code, a `default' branch was added, that
called \alg{int\_value()} with the node's subtype.
In an arbitrary context, it is possible that the syntax `tree' is not actually a tree, rather a cyclic graph, e.g., owing to data aliasing. Hence, there is a context in which the old function terminates whereas the new one is trapped in infinite recursion. \full{
Appendix~\ref{sec:sample}} \notfull{The full version of this article~\cite{EKS12}} includes the code of this function as well as an additional example in which mutual termination is not preserved.

In the remaining two comparisons \alg{rvt} marked correctly, in less than a minute each, that all mapped functions are mutually terminating.



\subsubsection{Conclusion and future research.}
We showed a proof rule for mutual termination, and a bottom-up decomposition algorithm for
handling whole programs. This algorithm calls a model-checker for discharging
the premise of the rule. Our prototype implementation of this algorithm in \alg{rvt} is the first to
give an automated (inherently incomplete) solution to the mutual termination
problem.

An urgent conclusion from our experiments is that checking mutual termination under free context is possibly insufficient, especially when it comes to programs that manipulate a global structure on the heap. Developers would also want to know whether their programs mutually terminate under the context of their specific program. \Long{This is not an easy modification to our algorithm, because the decomposition is based on a bottom-up traversal, hence ignoring the context. Perhaps a method can be found that propagates information down the call graphs that can restrict the context in which pairs of functions are checked for mutual termination. Another direction for future research is to improve the information that is propagated \emph{upwards}. Specifically we would like to refine the abstraction imposed by the use of uninterpreted functions. Adding \emph{function summaries} that provide more information about what these functions do -- making them more interpreted -- is bound to make the method more complete. A third direction } Another direction is to interface \alg{rvt} with an external tool that checks termination: in those cases that they can prove termination of one side but not of the other, we can use the results of Sec.~\ref{sec:term} to prove termination in the other side. We can also benefit from knowing that a pair of functions terminate (not just mutually terminate) because in such a case they should be excluded from the call-equivalence check of their callers.
%
Finally, it seems plausible to develop methods for proving termination by using the rule \mtermdp. One needs to find a variant of the input program that on the one hand is easier to prove terminating, and on the other hand is still call-equivalent to the original program.

%The limitations of the tool are inherited from \alg{rvt} itself, namely the fact that it does not support various features of C.


\bibliographystyle{abbrv}
\bibliography{biblio}
\appendix
\newpage
\section{Soundness proof for \mtermdp}
\label{sec:proof1}

For simplicity, assume that $\mapf$ provides a bijective mapping between the
functions in \comtag{m} (otherwise we can possibly satisfy this assumption via
inlining).
%
In the following by \emph{context} of a function call $f(\vec{in})$ we mean the execution from the time $f(\vec{in})$ entered the stack until it left it. We begin with two simple observations:
\begin{itemize}
\item[O1.] If there is a call $f(\vec{in})$ in the context of
    $f(\vec{in})$, then $f(\vec{in})$ does not terminate.

\item[O2.] If $f(\vec{in})$ terminates, then there can only be a
    finite number of parameter vectors in the context of $f(\vec{in})$ with which $f$ is called.
\end{itemize}
%
Denote by $\ufpdef{g}$ the {\emph set} (not a multiset) of parameter vectors
with which function $g$ is directly called from function $f$.

\begin{lemma}\label{lem:a1l1}
If $\requiv(\upbody{f},\upbody{f'}$), then $\requiv(f,f')$.
\end{lemma}

\begin{proof}: We prove the contrapositive of the lemma:
\[\lnot \requiv(f,f') \rightarrow \lnot \requiv(\upbody{f},\upbody{f'})\;.\]

Suppose \comandtag{f} are called with the same
value but for some \pairtag{g}$\in\mapf$ and input $\vec{in}$, $g(\vec{in})$ is called from $f$ and $g'(\vec{in})$ is not called from $f'$. This implies $\lnot \requiv(f,f')$. %$\vec{in}\in$\ufpdef{g}. %
In $\upbody{f}$, each call of $g$ is replaced with a call of
$\uf{g}$ with the same parameters. If the parameters values $\vec{in}$ do not depend on the result of any of the recursive functions, then the value of the parameters passed into the corresponding call
to $\uf{g}$ in $\upbody{f}$ will not depend on the result of any
uninterpreted function and, therefore, will be equal to the value of $\vec{in}$.
Otherwise, namely, $\vec{in}$ depends on the result of some (or several) recursive
function(s), then the non-deterministic values returned by the corresponding uninterpreted
function(s) will affect the parameter passed into the corresponding call
to $\uf{g}$ in $\upbody{f}$. But among all the possibilities for those
non-deterministic values, there are values that were actually returned by
the recursive functions replaced by those uninterpreted functions. Those
actual values caused that $\vec{in}$ was passed into the mentioned call of
$g$. Those actual values (returned by the uninterpreted functions) cause
$\uf{g}$ to be called with $\vec{in}$ in $\upbody{f}$. Hence,
$\vec{in}\in\ufp{g}$.

If there are no calls of $g'$ in $f'$ at all, then $\upbody{f'}$ contains
no calls of $\uf{g'}$ either. Hence, $\ufpp{g'}=\emptyset$, i.e.,
$\ufp{g}\neq\ufpp{g'}$. Otherwise (there is a call of $g'$ in $f'$)
by our assumption $\vec{in} \notin \ufpdefp{g'}$. Analogously, we can show
that in this case when all uninterpreted functions return actual values
that their corresponding original functions would, some value different
from $\vec{in}$ is passed into $\uf{g'}$ in $\upbody{f'}$. Hence, there
exists a computation with a set of non-deterministic values for which
$\vec{in}\notin\ufpp{g'}$, i.e., $\ufp{g}\neq\ufpp{g'}$.

So in both cases we found a computation where $\ufp{g}\neq\ufpp{g'}$, which implies $\lnot \requiv(\upbody{f},\upbody{f'})$. \qed \end{proof}

The next lemma proves the validity of the rule\mbox{\ \mtermd} for basic programs consisting of a single function. The function can be either simply recursive
\Long{as displayed in Fig.~\ref{fig:mostBasicRecCase}, }
or not recursive at all.
\Long{
\begin{figure}[here]
\begin{center}
\begin{minipage}{0.25\linewidth}
\begin{tikzpicture}
\tikzset{node distance = 1.35cm}
\Vertex[Math=true]{f}
\GraphLoop[dist=1.0cm,dir=NO,style={->, thick}](f)
\end{tikzpicture}
\end{minipage}
\begin{minipage}{0.25\linewidth}
\begin{tikzpicture}
\tikzset{node distance = 1.35cm}
\Vertex[Math=true]{f'}
\GraphLoop[dist=1.0cm,dir=NO,style={->, thick}](f')
\end{tikzpicture}
\end{minipage}
\end{center}
\caption{The most basic case of two recursive programs.}
\label{fig:mostBasicRecCase}
\end{figure}
}

\begin{lemma}\label{lem:a1l2}
If all the functions called in \pairtag{f}$\in\mapf$ are \comandtag{f} themselves, respectively, then $\requivpair{f} \rightarrow \mtpair{f}$.
\end{lemma}

\begin{proof}Consider \pairtag{f}$ \in\mapf$ called with the same parameter
$\vec{in}$ such that each one of \comandtag{f} can call only itself. Falsely assume $\lnot\mtpair{f}$. W.l.o.g., assume $\term{f(\vec{in})}$ and  $\lnot\term{f'(\vec{in})}$, where, recall, $\term{f(\vec{in})}$ denotes that $f(\vec{in})$ terminates.

Consider $\vec{in_1}$, the parameter passed into the non-returning
recursive call to $f'$. $\requivpair{f}$ implies that $f$ is called
with $\vec{in_1}$ too. $\vec{in_1} \neq \vec{in}$ because otherwise we would get $\lnot\term{f(\vec{in})}$ according to O1.

Now consider terminating $f(\vec{in_1})$ and non-terminating $f'(\vec{in_1})$.
The situation is quite similar to $f(\vec{in})$ and $f'(\vec{in})$. Thus we
conclude that $f(\vec{in_2})$ must call terminating $f(\vec{in_2})$, while
 $f(\vec{in_1})$ must call non-terminating $f'(\vec{in_2})$, such that $\vec{in_2}
\notin \{\vec{in}, \vec{in_1}\}$. We can go on descending the call stacks
discovering a new value $\vec{in_n} \notin \{\vec{in}, \vec{in_1}, ...,
\vec{in_{n-1}}\}$ with which both \comandtag{f} are called. But the number of
such unique values passed into $f$ is finite according to O2. Hence,
we must eventually reach a situation where either:
\begin{itemize}
\item $f(\vec{in_n})$ calls $f$ with some $\vec{in_i}$ that is already found
    in the call stack, which contradicts the assumption $\term{f(\vec{in})}$ (by O1);
\item $f(\vec{in_n)}$ calls $f$ with some value $\vec{in_{extra}}$ such that
    $f'(\vec{in_{extra}})$ is never called in $f'(\vec{in_n})$. It contradicts $\requivpair{f}$;
\item $f(\vec{in_n})$ has no more calls to $f$. However,  $f'(\vec{in_n})$ must call $f'$ because of $\lnot\term{f'(\vec{in_n})}$.
    Again, it contradicts $\requivpair{f}$.
\end{itemize}
Consequently, the assumption $\lnot\mtpair{f}$ cannot be true.
\qed
\end{proof}

%Denote the fact that function $f$ has a call to $g$ by $\calls{f}{g}$. Similarly, $\callspair{f}{g}$ denotes that $f$ has a call to $g$ and $f'$ has a call to $g'$.

The next lemma addresses mutual recursion with terminating outer calls.

\begin{lemma}\label{lem:a1l3}
The following rule is sound:
%For each function \pairtag{f} $\in\mapf(m)$ such that \pairtag{m}$\in\mapm$ the following:
\[
%\requivpair{f}\mbox{ }\land\mbox{ }\forall{\pairwtag{g}}.
%(\callspair{f}{g}) \rightarrow \left(
%\left(g\in C(m) \land g'\in C(m')\right) \rightarrow
%\pairwtag{g}\in\pairwtag{m} \lor
%\left(\term{g} \land \term{g'}\right)%\right)
\frac
{\begin{array}{l}
(\forall{\pairwtag{g}} \in \mapf.\ (\left(g \in C(m) \land g' \in C(m')\right) \rightarrow \left(\term{g} \land \term{g'}\right))\ \land \\
(\forall {\pairwtag{f}} \in\mapf(m).\ \requivpair{f})
\end{array}
}
{\forall {\pairwtag{f}} \in\mapf(m).\  \mtpair{f}}
\]
%if every function call in \pairtag{f}$\in\mapf(m)$ either terminates or calls to a function from \comtag{m},
% respectively, then $\requivpair{f}$
%implies $\mtpair{f}$.
\end{lemma}

\begin{proof}
(Proof sketch) Consider \pairtag{f}$\in$ \pairtag{m} that are called with the same parameter
$\vec{in}$. Assume by negation $\lnot\mtpair{f}$. W.l.o.g., assume $\term{f(\vec{in})}$
 and  $\lnot\term{f'(\vec{in})}$.

The premise of the lemma implies that all the outer calls (beyond \comtag{m}) terminate. Hence, only inner calls inside \comtag{m} could cause $\lnot\term{f'(\vec{in})}$.

From now on, the proof is very similar to that of Lemma~\ref{lem:a1l2}. We start traversing
the infinite call stack of $f'(\vec{in})$. The only difference is
that instead of descending with the same function pair in every call stack level
(i.e., \pairtag{f} in the proof of Lemma~\ref{lem:a1l2}), we now descend to some pair \pairtag{h_i} $\in$ \pairtag{m} in each level $\#i$.
According to O2, the number of unique values for
every function from $m$ found in the call stack of $f(\vec{in})$ is finite
because $\term{f(\vec{in})}$ is assumed. On the other side,
the number of functions in $m'$ is finite. Hence, the infinite call
stack of $f'(\vec{in})$ must include calls to some function $g'$ s.t. \pairtag{g} $\in$ \pairtag{m} which repeats an
infinite number of times. The latter will contradict either $\term{f(\vec{in})}$ or $\requivpair{g}$ (similarly to how the infinitely called $f'$ led to the final contradiction in the proof of Lemma~\ref{lem:a1l2}).
\qed
\end{proof}

\begin{lemma}\label{lem:a1l4}
The following rule is sound:
%For each function \pairtag{f} $\in\mapf(m)$ such that \pairtag{m}$\in\mapm$ the following:
\[
\frac
{\begin{array}{l}
(\forall{\pairwtag{g}} \in \mapf.\ (\left(g \in C(m) \land g' \in C(m')\right) \rightarrow \mtpair{g})\ \land \\
(\forall {\pairwtag{f}} \in\mapf(m).\ \requivpair{f})
\end{array}
}
{\forall {\pairwtag{f}} \in\mapf(m).\  \mtpair{f}}
\]
%\[
%\requivpair{f}\mbox{ }\land\mbox{ }\forall{\pairwtag{g}}. (\callspair{f}{g}) \rightarrow \left(\pairwtag{g}\in\pairwtag{m} \lor %\mtpair{g}\right)
%\]
%implies $\mtpair{f}$.
%$\requivpair{f}$ and for each function $g$ called by $f$ there exists a function $g'$ called by $f'$ (and vice-versa) such that \pairtag{g}$\in$\pairtag{m} $\lor\xspace{~}\mtpair{g}$, then $\mtpair{f}$.
\end{lemma}

\begin{proof}Consider \pairtag{f}$ \in\mapf$ called with the same parameter
$\vec{in}$. Assume by negation $\lnot\mtpair{f}$. W.l.o.g., assume $\term{f(\vec{in})}$
 and  $\lnot\term{f'(\vec{in})}$.

Consider any function call $g(\vec{in_1})$ in $f(\vec{in})$  s.t. $g \notin m \land \exists{g'}. $\pairtag{g}$\in\mapf$.
$\requivpair{f}$ implies that $f'(\vec{in_1})$ also calls $g'(\vec{in_1})$.
%$\calls{f'(\vec{in_1})}{g'(\vec{in_1})}$.
\pairtag{g}$\notin$\pairtag{m} implies $\mtpair{g}$, which implies that $g(\vec{in_1})$ and $g'(\vec{in_1})$ mutually terminate.%, g(\vec{in_1})}$.
%$\lnot\term{g'(\vec{in_1})} \leftrightarrow \lnot\term{g(\vec{in_1})}$.
Hence, $g'(\vec{in_1})$ must terminate because otherwise this would contradict $\term{f(\vec{in})}$.

Hence, all the outer function calls (referring beyond \comtag{m}) must terminate. Thus all the conditions satisfy the premise of Lemma~\ref{lem:a1l3}, which implies $\mtpair{f}$.
\qed
\end{proof}


\begin{theorem}
The inference rule\mbox{\ \mtermdp} is sound.
\end{theorem}
\begin{proof}
According to Lemma~\ref{lem:a1l1}, $\forall {\pairwtag{f}} \in\mapf(m).\ \requiv(\upbody{f},\upbody{f'})$ implies $\forall {\pairwtag{f}} \in\mapf(m).\ \requivpair{f}$. Thus, the second line of the premise of the rule of Lemma~\ref{lem:a1l4} is satisfied. The upper line of the premise in\mbox{\ \mtermdp} matches the upper line of the premise of the rule of Lemma~\ref{lem:a1l4}. Having all its premises satisfied, Lemma~\ref{lem:a1l4} implies: \[\forall {\pairwtag{f}} \in\mapf(m).\  \mtpair{f}\;.\]
\qed
\end{proof}


\comment{
\begin{lemma}\label{lem:a1l4}
For \pairtag{f}$\in\mapf(m)$ such that $m=\{f\}$ and $m'=\{f'\}$ the following holds: if any function called in \pairtag{f} is, respectively, either \comtag{f} itself, or some \comtag{g} such that \pairtag{g}$ \in\mapf\mbox{ } \land$ \pairtag{g}$\notin\mapf(m)$ and $\mtpair{g}$ is assumed, then $\requivpair{f}$ implies $\mtpair{f}$.
\end{lemma}

\begin{proof}Consider \pairtag{f}$ \in\mapf$ called with the same parameter
$\vec{in}$ such that each one of \comandtag{f} can call only itself. Assume by negation $\lnot\mtpair{f}$. W.l.o.g., assume $\term{f(\vec{in})}$
 and  $\lnot\term{f'(\vec{in})}$, where $\term{f(\vec{in})}$ denotes that $f(\vec{in})$ terminates.

Consider any function $g \neq f$ called in $f(\vec{in})$ with parameter
$\vec{X}$ such that \pairtag{g}$\in\mapf$. The call-equivalence of \pairtag{f}
implies that $f'(\vec{in})$ also calls $g'$ with the parameter $\vec{X}$. If
$g'(\vec{X})$ did not terminate, then $g(\vec{X})$ would not terminate either
because \pairtag{g} is mutually-terminating. Hence, no $g \neq f$ can prevent
the termination of $f'(\vec{in})$.

Consequently, only a recursive call to $f'$ itself can cause $f'(\vec{in})$ not
to terminate. Consider $\vec{X_1}$, the parameter passed into the non-returning
recursive call to $f'$. Since \pairtag{f} is call-equivalent, $f$ is called
with $\vec{X_1}$ too. $\vec{X_1} \neq \vec{in}$ because otherwise $f(\vec{in})$
would not terminate according to Lemma 1.

Now consider terminating $f(\vec{X_1})$ and non-terminating $f'(\vec{X_1})$.
The situation is quite similar to $f(\vec{in})$ and $f'(\vec{in})$. Thus we
conclude that there must be terminating $f(\vec{X_2})$ under $f(\vec{X_2})$ and
non-terminating $f'(\vec{X_2})$ under $f(\vec{X_1})$ such that $\vec{X_2}
\notin \{\vec{in}, \vec{X_1}\}$. We can go on descending the call stacks
discovering a new value $\vec{X_n} \notin \{\vec{in}, \vec{X_1}, ...,
\vec{X_{n-1}}\}$ with which both \comandtag{f} are called. But the number of
such unique values passed into $f$ is finite according to the Observation 2. Hence,
we must reach a situation where either:
\begin{itemize}
\item $f(\vec{X_n})$ calls $f$ with some $\vec{X_i}$ that is already found
    in the call stack, which contradicts the assumption that $f(\vec{in})$
    terminates (Observation 1);
\item $f(\vec{X_n)}$ calls $f$ with some value $\vec{Y}$ such that
    $f'(\vec{Y})$ is never called in $f'(\vec{X_n})$. It contradicts the
    call-equivalence of \pairtag{f};
\item $f(X_n)$ has no more calls to $f$. However, since $f'(\vec{X_n})$
    does not terminate and, as earlier shown, only infinite recursive calls
    to $f'$ can prevent its termination, $f'(\vec{X_n})$ must call $f'$.
    Again, it contradicts the call-equivalence of \pairtag{f}.
\end{itemize}
Consequently, the assumption that $f(\vec{in})$ and $f'(\vec{in})$ are not
mutually terminating cannot be true. \qed \end{proof}

Meanwhile, we have proven the soundness of the inference rule for simple
recursion only. The next lemma addresses mutual recursion.

\begin{lemma}\label{lem:a1l3}
For matched MSCC's \pairtag{m}, if all matched pairs of \pairtag{f} $\in$ \pairtag{m} have call-equivalent \pairbodies{f} and all the called functions which are not in \pairtag{m} belong to mutually-terminating pairs, then all the pairs of functions in \pairtag{m} are mutually-terminating.
\end{lemma}
\begin{proof} Proof sketch: according to Lemma~\ref{lem:a1l1}, if for all the pairs \pairtag{f} $\in$
\pairtag{m},   \pairbodies{f} is call-equivalent, then all those pairs are
call-equivalent by themselves.

Consider \pairtag{f}$\in$ \pairtag{m} called with the same parameter
$\vec{in}$. Similarly to the proof of Lemma~\ref{lem:a1l2}, we can show that only recursive
calls to some functions from \pairtag{m} can cause $f$ or $f'$ not to
terminate. Assume by negation \pairtag{f} is not mutually terminating. Without
loss of generality, assume that $f(\vec{in})$ terminates, while $f'(\vec{in})$
does not.

The number of functions in \pairtag{m} is finite. Hence, the infinite call
stack of $f'(\vec{in})$ must include calls to some function $g'$ that repeat an
infinite number of times such that \pairtag{g} $\in$ \pairtag{m}.

From now on, the proof is very similar to that of Lemma~\ref{lem:a1l2}. We start traversing
the infinite call stack of $f'(\vec{in})$ and compare every call to $h'$ to the
call of $h$ such that \pairtag{h} $\in$ \pairtag{m}. The only difference is
that instead of descending to the same function pair in every call stack level
(i.e., \pairtag{f} in the proof of Lemma~\ref{lem:a1l2}), we now descend to some pair among
the pairs of \pairtag{m}. According to Observation 2, the number of unique values for
every function from $m$ found in the call stack of $f(\vec{in})$ is finite
because $f(\vec{in})$ is assumed to terminate. On the other side, as mentioned
above, there exists infinitely repeating calls to $g'$ that will cause a
contradiction to the assumption that $f(\vec{in})$ terminates, while
$f'(\vec{in})$ does not: it can be shown that either $f(\vec{in})$ does not
terminate or \pairtag{g} are not call-equivalent. \qed \end{proof}

\begin{lemma}\label{lem:a1l4}
A call graph of a program is a directed acyclic graph of MSCCs (if we collapse the nodes of each MSCC into a single node).
\end{lemma}
\begin{proof}
Falsely assume that there is a loop between two MSCCs: $m_1$ and $m_2$. It implies existence of directed paths: $\pi_{1 \rightarrow 2}$ from $m_1$ to $m_2$ and $\pi_{2 \rightarrow 1}$ from $m_2$ to $m_1$. $\pi_{1 \rightarrow 2}$ implies that there is a call graph node $s_2$ in $m_2$ which is reachable from some node $t_1$ in $m_1$. $t_1$ is reachable from all the functions in $m_1$. $s_2$ can reach all the functions in $m_2$. Hence, every node in $m_1$ can reach each node in $m_2$. Analogously, the path $\pi_{2 \rightarrow 1}$ implies the reachability of each node in $m_1$ from every node of $m_2$. Hence, the set of the functions of $m_1 \cup m_2$ are strongly connected each to other. Since $m_1 \neq \emptyset \land m_2 \neq \emptyset \land m_1 \neq m_2$,  $m_1 \cup m_2$ is a bigger strongly connected component that includes each one of $m_1$ and $m_2$ in a contradiction to the fact that they are SCCs of maximal sizes.
\qed
\end{proof}

\begin{theorem}
The inference rule\mbox{\ \mtermdp} is sound.
\end{theorem}
\begin{proof}
Recall the MSCC  Denote the maximal distance of an MSCC node $m$ from any leaf
MSCC by $d(m)$ (If $m$ is a leaf MSCC node then $d(m) = 0$). The proof is by
induction on $d(m)$.

Base: Consider a pair of mapped leaf MSCC nodes: \pairtag{m} $\in \mapm$. $m$ has only calls to the functions of $m$ itself. $m'$ has only calls to the functions of $m'$ itself. By the assumption of the inference rule\mbox{\ \mtermd}, all matched pairs of \pairtag{f} $\in$ \pairtag{m} have call-equivalent \pairbodies{f}. Therefore, \pairtag{m} satisfies the premise of Lemma~\ref{lem:a1l3}. Consequently, the functions of this pair are mutually terminating.

Step: Assume that using the rule each mapped pair of MSSC nodes \pairtag{t} such that $max\{d(t), d(t')\} \le i$ was established as mutually terminating. We must prove that the rule is true for any pair of mapped MSSC nodes \pairtag{m} $\in \mapm$ such that $max\{d(m), d(m')\} = i+1$.

Consider a pair of mapped leaf MSCC nodes, \pairtag{m} $\in \mapm$ such that $max\{d(m), d(m')\} = i+1$.
If $m$ had an outer function call to an MSCC node $x$ such that $d(x) \ge i+1$, it would contradict the assumption that $d(m) \le i+1$. Hence, every outer call from $m$ refers to a function from an MSCC node $t$ for which $d(t) \le i$. The assumption of the induction implies that the functions are established as mutually terminating. Analogously, all the outer calls in $m'$ refer to functions established as mutually terminating. The assumption of the inference rule\mbox{\ \mtermd} implies that all matched pairs of \pairtag{f} $\in$ \pairtag{m} have call-equivalent \pairbodies{f}. Hence, \pairtag{m} satisfies the premise of Lemma~\ref{lem:a1l3}. Consequently, the functions of this pair are mutually terminating.
\qed
\end{proof}
}

\full{
\section{Sample code}
\label{sec:sample}
The following two code snippets from the \alg{Betik} project demonstrate changes that do not preserve mutual termination. Explanations are found in Sec.~\ref{sec:experiments}.

\begin{figure}[top]
\begin{tabular}{p{5 cm}|p{7 cm}} \hline
\begin{algorithmic}
%\Function{$int\_value$}{value\_t *$v$}
\State value\_t *\alg{int\_value}(value\_t *$v$) \{
\State \hspace{0.35cm}{\bf switch}($v$$\rightarrow$$type$) \{
\State \hspace{0.35cm}\hspace{0.35cm}{\bf case} 0: $v = \ldots$; {\bf break};
\State \hspace{0.35cm}\hspace{0.35cm}$\vdots$
\State \hspace{0.35cm}\hspace{0.35cm}{\bf case} N: $v = \ldots$; {\bf break};
\State \hspace{0.35cm}\}
\State \hspace{0.35cm}\Return $v$;
\State \}
%\EndFunction
\end{algorithmic}
&
\begin{algorithmic}
%\Function{int\_value'}{value\_t *$v'$}
\State value\_t *\alg{int\_value'}(value\_t *$v'$) \{
\State \hspace{0.35cm}{\bf switch}($v'$$\rightarrow$$type$) \{
\State \hspace{0.35cm}\hspace{0.35cm}{\bf case} 0: $v' = \ldots$; {\bf break};
\State \hspace{0.35cm}\hspace{0.35cm}$\vdots$
\State \hspace{0.35cm}\hspace{0.35cm}{\bf case} N: $v' = \ldots$; {\bf break};
\State \hspace{0.35cm}\hspace{0.35cm}{\bf default}: $v'$ = \alg{int\_value'}($v'$$\rightarrow$$subvalue$);
%\{val\_t *$u'=\ldots$; $int\_value'$($u'$);\}
\State \hspace{0.35cm}\}
\State \hspace{0.35cm}\Return$v'$;
\State \}
%\EndFunction
\end{algorithmic}
\\
\hline
\end{tabular}
\caption{Two possibly non-mutually terminating versions of \alg{int\_value()}.}\label{fig:int_value}
\end{figure}


The code in Fig.~\ref{fig:parse_funccall} is a function called \alg{parse\_funccall} that receives a pointer to a function node and processes it according to the function name. The newer version handles an additional option for the function name and calls a new function \alg{list\_set\_item}. The latter receives a pointer to a list, traverses it from the list head, and modifies data of some of its items. The traversal ends upon reaching a `NULL' node, but may not terminate in an arbitrary context, e.g., when the list is cyclic. The new function is not mapped to any function in the old code. Indeed Alg.~\ref{alg:ProveMT} aborts at line~\ref{step:generate_map} when encountering this function.

\begin{figure}
\begin{tabular}{p{6.1 cm}|p{6.1 cm}} \hline
\begin{algorithmic}
\Function{parse\_funccall}{fcall\_t *$f$} \{
%\State variable\_t *$int\_funccall$(..., funccall\_t* $f$) \{
\State {\bf if} (!$strcmp$($f$$\rightarrow$func\_name, ``env")) \{
\State \hspace{0.4cm}$\ldots$
\State \} {\bf else if} (!$strcmp$($\ldots$, ``len")) \{
\State \hspace{0.4cm}$\ldots$
\State \}
\EndFunction
\State \}
\end{algorithmic}
&
\begin{algorithmic}
\Function{parse\_funccall'}{fcall\_t *$f'$} \{
%\State variable\_t *$int\_funccall'$(..., funccall\_t* $f'$) \{
\State {\bf if} (!$strcmp$($f'$$\rightarrow$func\_name, ``env")) \{
\State \hspace{0.4cm}$\ldots$
\State \} {\bf else if} (!$strcmp$($\ldots$, ``len")) \{
\State \hspace{0.4cm}$\ldots$
\State \} {\bf else if} (!$strcmp$($\ldots$, ``set")) \{
\State \hspace{0.4cm}list\_t *$list' = \ldots$;
\State \hspace{0.4cm}\alg{list\_set\_item'}($list'$, $\ldots$);
\State \}
\EndFunction
\State \}
%\comment{
  \\
  \Function{list\_set\_item'}{list\_t *$list'$, ...} \{
  \State listitem\_t *$item' = list'$$\rightarrow$head;
  \State {\bf while}($item'$ != NULL) \{
  \State \hspace{0.4cm}$\ldots$
  \State \hspace{0.4cm}$item' = item'$$\rightarrow$next;
  \State \}
  \EndFunction
  \State \}
%}
\end{algorithmic}
\\
\hline
\end{tabular}
\caption{Two possibly non-mutually terminating versions of \alg{parse\_funccall()} and a newly introduced non-mapped function \alg{list\_set\_item'()}.}
\label{fig:parse_funccall}
\end{figure}
}


%\subsection{Soundness proof for \termd} \label{sec:proof2}
%The proof follows similar lines to that of \mtermdp. We give a proof sketch.
%Falsely assume that there is a function $f'$ in $m'$ that does not terminate, whereas for all $\pair{g,g'} \in \mapf(m)$, $\callc(g,g')$.
%There exists a value $in$ such that $f'(in)$ does not terminate. The infinite call stack of $f'(in)$ must contain a call, say from $h'(in_1)$ to $g'(in_2)$, whereas $h(in_1)$ does not call $g(in_2)$ in the call stack of $f(in)$ (assuming $\{$\pairtag{g}, \pairtag{h}$\} \subseteq \mapf$).
%This contradicts our premise that $\callc(h,h')$.

\full{
\section{Mutual termination of functions with infinite-type variables} \label{sec:infinite}
The question whether call equivalence can be proven even when the variable types are infinite is particularly interesting in cases such as the Collatz problem mention in Example~\ref{ex:collatz}, because if the answer is yes, then this stands in stark contrast to the corresponding termination problem (recall that termination of this program is not known). Indeed in this case call equivalence can be verified even when the input $a$ is a true (infinite) integer. Rather than the model-theoretic solution described so far (which relies on CBMC's ability to reason about C programs, in which variables are of a finite type), here we take the proof-theoretic approach instead and formulate a verification condition which is valid only if the two functions are call-equivalent. For this purpose we need to represent the transition relation of the two functions $T_{\upbody{f}}, T_{\upbody{f'}}$, which is easy to do with the help of static single assignment form:


\begin{tabular}{p{6 cm}p{6 cm}} \hline
\begin{algorithmic}
\Function {$f^{UF}$}{{\bf int} $a_0$}
\State {\bf int} $even_0 := 0, ret_0 := 0,$\\  \mbox{ }$\qquad a_1, ret_1, ret_2, even_1, even_2$;
\If {$a_0 > 1$}
\If {$\lnot(a_0\ \%\ 2)$} %\Comment{even}
    \State $a_1 := a_0/2$;
    \State $even_1 := 1$;
    \Else \ $a_1 := 3a_0 + 1$;
\EndIf
\State $even_2 := \lnot(a_0\ \%\ 2)\ ?\ even_1 : even_0$;
\State $ret_1 := even_2$ + \alg{UF}($f,a_1$);
\EndIf
\State $ret_2 := a_0 > 1\ ?\ ret_1 : ret_0$;  
\State \Return $ret_2$;
\EndFunction
\end{algorithmic}
&
\begin{algorithmic}
\Function {$f'^{UF}$}{{\bf int} $a_0'$}
\State {\bf int} $t_0', odd_0' := 0, ret_0' := 0,$\\ \mbox{ }$\qquad a_1', t_1', odd_1', odd_2', ret_1'$;
\If {$a_0' \leq 1$} \Return $ret_0'$; \EndIf
\State 	$t_1' := a_0' / 2$;
\If {$a_0'\ \%\ 2$}  % \Comment{odd}
\State		$a_1' := 6 t_1' + 4$;
\State		$odd_1' := 1$;
\Else \ $a_1' := t_1'$;	
\EndIf
\State $odd_2' := a_0' \% 2\ ?\ odd_1' : odd_0'$;
\State $ret_1' := odd_2'$ + \alg{UF}($f',a_1'$);
\State \Return $ret_1'$;
\EndFunction
\end{algorithmic} \\ \hline
\end{tabular}
\vspace{0.3 cm}



\comment{
\vspace{0.3 cm}
\noindent
\begin{tabular}{p{6 cm}p{6 cm}} \hline
\begin{algorithmic}
\Function {$f^{UF}$}{int $a_0$}

\If{$a_0 > 1$}

\If {$\lnot$($a_0 \% 2$)}

\State $a_1 := a_0/2$;

\Else\ $a_1 := 3a_0 + 1$;

\EndIf

\State $\uf{f}(a_1)$;

\State \Return 0;
\EndIf
\EndFunction
\end{algorithmic}
&
\begin{algorithmic}
\Function {$f'^{UF}$}{int $a'_0$}

\If{$a'_0 \leq 1$} \Return 1; \EndIf

\State $t'_0 := a'_0  / 2$;

\If {$a'_0\% 2$}

\State $a'_1 := 6t'_0 + 4$;

\Else\ $a'_1 := t'_0$;

\EndIf

\State $\uf{f'}(a'_1)$;

\State \Return 1;
\EndFunction
\end{algorithmic} \\ \hline
\end{tabular}
\vspace{0.3 cm}
}

\noindent It is not hard to see that the following verification condition is valid if $f,f'$ are call-equivalent, and can be decided with a decision procedure for integer linear arithmetic\footnote{The only operator which is not explicitly linear here is the modulo ($\%$) sign, but it can be reduced to integer linear arithmetic based on the equivalence $x \% 2 = y \leftrightarrow 2 * z + y = x$ for a new integer variable $z$.}.
%
\begin{equation}
(T_{\upbody{f}} \ \land\ T_{\upbody{f'}}\ \land\ a_0 = a'_0)\ \rightarrow
\big(((a_0 > 1) \leftrightarrow \lnot(a'_0 \leq 1))\ \land\ ((a_0 > 1) \rightarrow (a_1 = a'_1)) \big)\;.
\end{equation}
%


} 
%%%%%%%%%%%%%%%%%%%%%%%% A decomposition algorithm %%%%%%%%%%%%%%%%%%%%%%


\end{document}
